\documentclass[11pt,a4paper]{article}
\usepackage{amsmath,amssymb,amsfonts,mathtools}
\usepackage{bm}
\usepackage{booktabs}
\usepackage{array}
\usepackage{graphicx}
\usepackage{geometry}
\geometry{margin=2.4cm}
\usepackage{fontspec}
\setmainfont{CMU Serif}
\setsansfont{CMU Sans Serif}
\setmonofont{CMU Typewriter Text}
\usepackage[english,russian]{babel}
\usepackage[hidelinks]{hyperref}
\usepackage{enumitem}
\providecommand{\dd}{\mathrm{d}}
\title{\textbf{GreenTensor}: аналитические решатели рассеяния на основе\\ тензорной функции Грина --- план и математика по геометриям}
\author{Проект GreenTensor (УрФУ)}
\date{\today}
\begin{document}
\maketitle
\begin{abstract}
\noindent Документ собирает единый план и полную математику аналитических (на основе диадной/тензорной функции Грина и разделения переменных) решателей задачи электромагнитного рассеяния для строгого семейства конечных примитивов --- сферы, вытянутого и сплюснутого сфероида и трёхосного эллипсоида --- с поддержкой многослойных сред и объединением в сложную геометрию методом суперпозиции $T$-матриц (GMM). Сфера является каноническим математическим ядром. Каждый раздел сопровождается ссылками на публикации, против которых верифицируются формулы.
\end{abstract}
\tableofcontents
\bigskip

\section{Постановка и архитектура решения}
\label{plan-sec:plan}

\subsection{Цель проекта и круг решаемых задач}

Рассматривается классическая задача дифракции монохроматической электромагнитной волны на конечном диэлектрическом (в общем случае поглощающем, магнитодиэлектрическом и многослойном) теле, помещённом в однородную внешнюю среду. Цель проекта \textsc{GreenTensor} --- построение \emph{строго аналитических} решателей этой задачи на основе метода диадной функции Грина и разделения переменных, без обращения к объёмной или поверхностной дискретизации. В отличие от сеточных подходов (FEM, FDTD, MoM), аналитическое решение даёт замкнутые выражения для коэффициентов рассеяния, контролируемую сходимость ряда по мультиполям и машинную точность при умеренной вычислительной стоимости, что и делает его эталоном (арбитром) для верификации численных пакетов \cite{bohren1983,mishchenko2002}.

Принята временн\'{а}я конвенция $e^{-i\omega t}$; расходящаяся (исходящая) волна описывается сферической функцией Ханкеля первого рода $h_n^{(1)}$, что в терминах функций Риккати--Бесселя соответствует $\xi_n=\psi_n-i\chi_n$. Все среды считаются линейными, изотропными и локальными, с комплексными $\varepsilon$ и $\mu$; поглощению отвечает $\operatorname{Im}\varepsilon>0$.

\subsection{Управляющие уравнения и операторная (интегральная) форма}

Вне источников поле $\bm{E}$ удовлетворяет векторному уравнению Гельмгольца
\begin{equation}
\nabla\times\nabla\times\bm{E}(\bm{r})-k^2(\bm{r})\,\bm{E}(\bm{r})=\bm{0},
\qquad k^2(\bm{r})=\omega^2\varepsilon(\bm{r})\mu(\bm{r}),
\label{plan-eq:vecthelm}
\end{equation}
с условиями непрерывности тангенциальных компонент $\bm{E}$ и $\bm{H}$ на каждой границе раздела и условием излучения Зоммерфельда на бесконечности \cite{stratton1941}. Полное поле раскладывается на падающее и рассеянное, $\bm{E}=\bm{E}^{\mathrm{inc}}+\bm{E}^{\mathrm{sca}}$, причём $\bm{E}^{\mathrm{inc}}$ регулярно во всём пространстве, а $\bm{E}^{\mathrm{sca}}$ удовлетворяет условию излучения.

Эквивалентно \eqref{plan-eq:vecthelm} может быть записано в виде объёмного интегрального уравнения Липпмана--Швингера. Вводя фоновую (свободную) диадную функцию Грина $\bar{\bm{G}}_0(\bm{r},\bm{r}')$ внешней среды, удовлетворяющую
\begin{equation}
\nabla\times\nabla\times\bar{\bm{G}}_0-k_{\mathrm{ext}}^2\,\bar{\bm{G}}_0
=\bar{\bm{I}}\,\delta(\bm{r}-\bm{r}'),
\qquad k_{\mathrm{ext}}^2=\omega^2\varepsilon_{\mathrm{ext}}\mu_{\mathrm{ext}},
\label{plan-eq:greendyad}
\end{equation}
полное поле выражается через объёмный интеграл по телу $V$ с контрастом волнового числа в качестве источника:
\begin{equation}
\bm{E}(\bm{r})=\bm{E}^{\mathrm{inc}}(\bm{r})
+\int_{V}\bar{\bm{G}}_0(\bm{r},\bm{r}')\cdot
\big[k^2(\bm{r}')-k_{\mathrm{ext}}^2\big]\,\bm{E}(\bm{r}')\,\mathrm{d}^3 r'.
\label{plan-eq:lippmann}
\end{equation}
Для немагнитной среды ($\mu\equiv\mu_0$) контраст сводится к $k^2-k_{\mathrm{ext}}^2=\omega^2\mu_0\,\Delta\varepsilon$ с диэлектрическим контрастом $\Delta\varepsilon(\bm{r})=\varepsilon(\bm{r})-\varepsilon_{\mathrm{ext}}$. Уравнение \eqref{plan-eq:lippmann} --- операторная постановка рассеяния. Аналитические решатели проекта не решают \eqref{plan-eq:lippmann} напрямую, а используют его как концептуальную основу: внешнее рассеянное поле полностью определяется линейным откликом тела на падающее поле, что и формализуется $T$-матрицей (см.~\S\ref{plan-ssec:tmatrix}).

\subsection{Строгое семейство аналитических примитивов и критерий отбора}

Векторное уравнение \eqref{plan-eq:vecthelm} (а вместе с ним и скалярное уравнение Гельмгольца, к которому сводятся потенциалы Дебая) допускает разделение переменных лишь в специальном классе криволинейных координат. Согласно классификации Эйзенхарта (условия Штеккеля), скалярное уравнение Гельмгольца разделяется ровно в \emph{одиннадцати} ортогональных системах координат конфокальных квадрик и их вырождений \cite{eisenhart1934}. Для задачи о \emph{конечном} рассеивателе релевантны лишь те системы, в которых граница тела совпадает с \emph{одной} координатной поверхностью семейства (тогда граничные условия не перемешивают гармоники базиса). Этому требованию отвечают:
\begin{itemize}
  \item \textbf{сфера} --- координатная поверхность $r=\mathrm{const}$ в сферической системе;
  \item \textbf{вытянутый и сплюснутый сфероид} (prolate / oblate) --- поверхность $\xi=\mathrm{const}$ в вытянутых/сплюснутых сфероидальных координатах;
  \item \textbf{трёхосный эллипсоид} --- поверхность $\xi=\mathrm{const}$ в эллипсоидальных (Лам\'{е}) координатах.
\end{itemize}
Эти три тела и образуют строгое семейство примитивов проекта. Прочие из одиннадцати систем (параболоид, цилиндры, конус и т.\,д.) либо описывают бесконечные/незамкнутые поверхности, либо не дают конечного тела, ограниченного единственной координатной поверхностью, и потому исключаются. Сфера выступает каноническим математическим ядром: для неё разделение полное, $T$-матрица диагональна (коэффициенты Ми), а реализация (\texttt{green\_tensor/01\_sphere.py}, \texttt{green\_tensor/mie\_core.py}) служит эталоном для остальных примитивов.

\subsection{Объединяющий принцип: общий базис ВСВФ и \texorpdfstring{$T$}{T}-матрица}
\label{plan-ssec:tmatrix}

Падающее и рассеянное поля разлагаются по векторным сферическим волновым функциям (ВСВФ) $\bm{M}_{n}^{(J)},\bm{N}_{n}^{(J)}$, где $J=1$ отвечает регулярным (бесселевым $j_n$), $J=3$ --- исходящим (ханкелевым $h_n^{(1)}$) решениям \cite{stratton1941,mishchenko2002}. Обозначив мультииндекс $n=(l,m)$,
\begin{align}
\bm{E}^{\mathrm{inc}}(\bm{r}) &=\sum_{n}\Big[d_n^{M}\,\bm{M}_n^{(1)}(\bm{r})+d_n^{N}\,\bm{N}_n^{(1)}(\bm{r})\Big],
\label{plan-eq:incexp}\\
\bm{E}^{\mathrm{sca}}(\bm{r}) &=\sum_{n}\Big[c_n^{M}\,\bm{M}_n^{(3)}(\bm{r})+c_n^{N}\,\bm{N}_n^{(3)}(\bm{r})\Big].
\label{plan-eq:scaexp}
\end{align}
В силу линейности задачи \eqref{plan-eq:vecthelm} связь коэффициентов разложения --- линейная и задаётся \emph{матрицей перехода} (transition matrix, $T$-матрицей) \cite{waterman1971}:
\begin{equation}
\bm{c}=\bm{T}\,\bm{d},
\qquad
\begin{pmatrix}\bm{c}^{M}\\[2pt]\bm{c}^{N}\end{pmatrix}
=\begin{pmatrix}\bm{T}^{11}&\bm{T}^{12}\\[2pt]\bm{T}^{21}&\bm{T}^{22}\end{pmatrix}
\begin{pmatrix}\bm{d}^{M}\\[2pt]\bm{d}^{N}\end{pmatrix}.
\label{plan-eq:Tdef}
\end{equation}
$T$-матрица \eqref{plan-eq:Tdef} зависит только от тела (геометрии и материала), но не от направления и поляризации падающей волны, что и делает её удобным единым описанием рассеивателя \cite{waterman1971,mishchenko2002}. Объединяющий принцип проекта формулируется так: \emph{каждый примитив поставляет свою $T$-матрицу в едином базисе сферических ВСВФ}. Для сферы $\bm{T}=\operatorname{diag}(-a_n,-b_n)$ диагональна (коэффициенты Ми; $m$ и тип мультиполя не связаны), причём знак отвечает соглашению $\bm{S}=\bm{I}+2\bm{T}$ для матрицы рассеяния \cite{waterman1971,bohren1983}. Для сфероида естественный базис --- сфероидальные ВСВФ; полученная в нём блочная (по $m$) $T$-матрица \emph{перепроецируется} на сферический базис, после чего она становится совместимой со сферой и эллипсоидом.

Сложная геометрия собирается из примитивов методом обобщённой матрицы Ми (Generalized Multiparticle Mie, GMM): поля, исходящие из одного тела, выражаются как падающие для другого с помощью \emph{трансляционных теорем сложения} для ВСВФ \cite{stein1961,cruzan1962}. Если $\bm{T}_p$ --- $T$-матрица $p$-го тела, а $\bm{\alpha}(\bm{r}_{pq})$ --- оператор трансляции из центра $q$ в центр $p$, то самосогласованная система для коэффициентов возбуждения $\bm{a}_p$ имеет вид
\begin{equation}
\bm{a}_p=\bm{d}_p+\sum_{q\neq p}\bm{\alpha}(\bm{r}_{pq})\,\bm{T}_q\,\bm{a}_q,
\qquad
\bm{c}_p=\bm{T}_p\,\bm{a}_p,
\label{plan-eq:gmm}
\end{equation}
где $\bm{d}_p$ --- разложение внешней падающей волны относительно центра $p$ \cite{xu1995}. Система \eqref{plan-eq:gmm} замыкает многочастичную задачу через те же сферические ВСВФ, что и одиночный примитив, --- отсюда требование общего базиса.

\subsection{Многослойность и условие конфокальности}

Все примитивы должны поддерживать \emph{стратифицированную} (многослойную) структуру. Способ построения многослойного решения и его область применимости радикально различаются по геометриям.

\paragraph{Сфера.} Концентрические слои допускают \emph{произвольные} радиусы границ. Разделение остаётся полным, азимутальный порядок $m$ и тип (TE/TM) не связываются, а $T$-матрица остаётся диагональной. Отклик каждой моды получается скалярной рекурсией импеданса (TM) / адмитанса (TE) от ядра к внешней границе --- рекурсией Адена--Керкера \cite{aden1951,bohren1983}. В реализации (\texttt{mie\_core.py}) это рекурсия по $Z_n,Y_n$ от слоя $h-1$ к слою $h$:
\begin{equation}
Z_n^{(h)}=\sqrt{\tfrac{\varepsilon_{h+1}}{\varepsilon_h}}\;
\frac{C_n'+Z_n^{(h-1)}\,S_n'}{C_n+Z_n^{(h-1)}\,S_n},
\label{plan-eq:adenkerker}
\end{equation}
где $C_n,S_n$ и их производные $C_n',S_n'$ --- кросс-произведения функций Риккати--Бесселя $\psi_n,\chi_n$ соседних слоёв (например, $C_n=\psi_n(k_h a_h)\,\chi_n'(k_h a_{h-1})-\chi_n(k_h a_h)\,\psi_n'(k_h a_{h-1})$), а $\sqrt{\varepsilon_{h+1}/\varepsilon_h}$ --- скачок импеданса на границе.

\paragraph{Сфероид.} Аналитическое многослойное решение существует \emph{только для конфокальных оболочек}: все границы имеют общее (полу)фокусное расстояние, т.\,е. $a_j^2-b_j^2=(d/2)^2$ для каждого слоя $j$, где $d$ --- межфокусное расстояние, а $d/2$ --- полуфокусное \cite{farafonov1996,farafonov2012}. Следствие: соотношение полуосей (aspect ratio) меняется от слоя к слою и не может задаваться независимо. В отличие от сферы, разделение неполно: при фиксированном $m$ мультипольные порядки $l$ связываются, поэтому рекурсия --- \emph{блочно-матричная} по каждому азимутальному $m$ (порядки $m$ между собой не связаны). Расщепление на TE/TM имеет место лишь при осевом падении ($m=0$) \cite{asano1975,farafonov2012}.

\paragraph{Трёхосный эллипсоид.} Полноволновая векторная задача \emph{не разделяется}: волновые функции Лам\'{е} практически не вычислимы за пределами низших степеней. Строгое аналитическое многослойное решение существует \emph{только в квазистатическом (рэлеевском) пределе}: конфокальные оболочки, уравнение Лапласа, передача через $2\times2$ матрицу факторов деполяризации (формализм Сиволы--Линделла) \cite{sihvola1990}.

\subsection{Сводная таблица применимости}

\begin{table}[ht]
\centering
\begin{tabular}{@{}llll@{}}
\toprule
Геометрия & Однородное тело & Многослойность & Базис / метод \\
\midrule
Сфера & полноволновое (Ми) & произвольные радиусы; диаг.\ $T$ & сферич.\ ВСВФ; рек.\ Адена--Керкера \\
Сфероид (prolate/oblate) & полноволновое & только конфокальные оболочки & сфероид.\ ВСВФ; блок по $m$ \\
Трёхосный эллипсоид & квазистатика (Рэлей) & только конфокальные, квазистатика & потенциал Лапласа; $2\times2$ деполяр. \\
\bottomrule
\end{tabular}
\caption{Применимость строгого аналитического решения по геометриям и режимам. Полноволновая векторная задача разделяется для сферы и сфероида; для трёхосного эллипсоида строгое решение доступно лишь квазистатически.}
\label{plan-tab:feasibility}
\end{table}

\subsection{Стратегия верификации}

Корректность решателей подтверждается многоуровнево. (1)~\emph{Внутренняя согласованность}: проверка баланса энергии $Q_{\mathrm{ext}}=Q_{\mathrm{sca}}+Q_{\mathrm{abs}}$ (оптическая теорема) и унитарности $T$-матрицы для непоглощающих тел, $\bm{T}^{\dagger}\bm{T}+\operatorname{Re}\bm{T}=\bm{0}$ \cite{waterman1971,mishchenko2002}. (2)~\emph{Эталон сферы}: коэффициенты Ми сравниваются с независимой аналитической реализацией (арбитром), включая поглощающие и многослойные конфигурации \cite{bohren1983}. (3)~\emph{Предельные переходы}: сфероид при стремлении эксцентриситета к нулю должен воспроизводить сферу; квазистатический эллипсоид --- известные поляризуемости через факторы деполяризации \cite{sihvola1990,bohren1983}. (4)~\emph{Кросс-валидация} с независимыми численными пакетами (FEM/MoM) на сечениях рассеяния и диаграммах направленности.

Дальнее поле и сечения вычисляются из коэффициентов \eqref{plan-eq:scaexp} стандартным образом; для сферы (нормировка на $\pi R^2$)
\begin{equation}
Q_{\mathrm{sca}}=\frac{2}{x^2}\sum_{n=1}^{\infty}(2n+1)\big(|a_n|^2+|b_n|^2\big),\qquad
Q_{\mathrm{ext}}=\frac{2}{x^2}\sum_{n=1}^{\infty}(2n+1)\operatorname{Re}(a_n+b_n),
\label{plan-eq:qsca}
\end{equation}
где $x=k_{\mathrm{ext}}R$ --- параметр размера \cite{bohren1983}. Соотношения \eqref{plan-eq:qsca} реализованы в ядре и служат первой контрольной точкой всей цепочки.

\section{Слоистая сфера: каноническое ядро}
\label{sphere-sec:sphere}

Сфера --- единственное тело из рассматриваемого семейства, для которого векторная задача рассеяния решается аналитически в \emph{замкнутом} виде при произвольном числе концентрических слоёв. Поэтому радиально-слоистая сфера служит каноническим математическим ядром всего комплекса ТФГ: T-матрица здесь диагональна, рекурсия скалярна (поканальна), а полученные коэффициенты $M_n,N_n$ напрямую отождествляются (с точностью до фазовой конвенции) с коэффициентами Ми $a_n,b_n$. В этом разделе мы проводим вывод от уравнений Максвелла до сечений рассеяния, согласуя все обозначения с верифицированной реализацией (\texttt{green\_tensor/01\_sphere.py}, \texttt{green\_tensor/mie\_core.py}).

\subsection{От уравнений Максвелла к векторному уравнению Гельмгольца}

Принята временн\'ая конвенция $e^{-i\omega t}$. В кусочно-однородной среде без сторонних токов уравнения Максвелла для гармонических полей имеют вид \cite{stratton1941,bohren1983}
\begin{align}
\nabla\times\bm{E} &= i\omega\mu\,\bm{H}, &
\nabla\times\bm{H} &= -i\omega\varepsilon\,\bm{E},\\
\nabla\cdot(\varepsilon\bm{E}) &= 0, &
\nabla\cdot(\mu\bm{H}) &= 0.
\end{align}
Внутри каждого слоя $j$ с постоянными $(\varepsilon_j,\mu_j)$ исключение $\bm{H}$ даёт векторное уравнение Гельмгольца
\begin{equation}
\nabla\times\nabla\times\bm{E}-k_j^2\bm{E}=\bm{0},
\qquad
k_j=\omega\sqrt{\varepsilon_j\mu_j}=k_0\,\sigma_j,
\quad
\sigma_j=\sqrt{\varepsilon_j\mu_j},
\label{sphere-eq:helmholtz}
\end{equation}
где $k_0=\omega/c$ --- волновое число вакуума, а $\sigma_j$ --- комплексный показатель преломления слоя. В коде это $\sigma_j=\texttt{sqrt(eps*miy)}$, причём поглощение ($\operatorname{Im}\varepsilon_j>0$) входит \emph{в аргумент} специальных функций через комплексное $k_j$, а не как модуль; это согласуется с принятой конвенцией $e^{-i\omega t}$ \cite{bohren1983}.

\subsection{Векторные сферические волновые функции}

Решения \eqref{sphere-eq:helmholtz} строятся из скалярной порождающей функции $\psi_{mn}$, удовлетворяющей скалярному уравнению Гельмгольца $\nabla^2\psi+k^2\psi=0$, которое в сферических координатах разделяется \cite{eisenhart1934,stratton1941}:
\begin{equation}
\psi_{mn}=z_n(kr)\,P_n^{m}(\cos\theta)\,e^{im\varphi},
\label{sphere-eq:scalar-gen}
\end{equation}
где $z_n$ --- одна из сферических функций Бесселя, $P_n^m$ --- присоединённая функция Лежандра. Векторные сферические волновые функции (VSWF) определяются через $\psi_{mn}$ и вектор положения $\bm{r}$ \cite{stratton1941,bohren1983,mishchenko2002}:
\begin{align}
\bm{M}_{mn} &= \nabla\times(\bm{r}\,\psi_{mn}), &
\bm{N}_{mn} &= \frac{1}{k}\,\nabla\times\bm{M}_{mn}.
\label{sphere-eq:MN-def}
\end{align}
Поле $\bm{E}$ (и через $\nabla\times$ --- поле $\bm{H}$) раскладывается в ряд по $\bm{M}_{mn},\bm{N}_{mn}$. Здесь $\bm{M}_{mn}$ чисто поперечны ($\bm{r}\cdot\bm{M}_{mn}=0$, моды типа TE относительно радиуса), а $\bm{N}_{mn}$ несут радиальную компоненту (моды типа TM). При осевом падении ($e^{im\varphi}$ с $m=1$ для плоской волны) азимутальная зависимость фиксирована, и задача сводится к одномерным радиальным уравнениям для каждого порядка $n$.

\subsection{Радиальные функции Риккати--Бесселя}

Радиальная зависимость удобно записывается через функции Риккати--Бесселя \cite{bohren1983,dlmf}:
\begin{align}
\psi_n(x) &= x\,j_n(x), &
\chi_n(x) &= -x\,y_n(x), &
\xi_n(x) &= x\,h_n^{(1)}(x)=\psi_n(x)-i\,\chi_n(x),
\label{sphere-eq:riccati}
\end{align}
где $j_n,y_n,h_n^{(1)}=j_n+i\,y_n$ --- сферические функции Бесселя, Неймана и Ханкеля первого рода; знак в $\chi_n=-xy_n$ соответствует конвенции \cite{bohren1983}, при которой $\xi_n=\psi_n-i\chi_n=x(j_n+iy_n)$. Выбор $h_n^{(1)}$ для рассеянного поля обусловлен конвенцией $e^{-i\omega t}$: $h_n^{(1)}(x)\sim(-i)^{n+1}e^{ix}/x$ при $x\to\infty$ задаёт \emph{исходящую} волну \cite{bohren1983,mishchenko2002}. В реализации $\xi_n$ вычисляется напрямую через \texttt{hankel1} (функция \texttt{\_xi}), а сингулярное решение хранится без знака как $x\,y_n(x)$; этот выбор глобального знака $\chi_n$ \emph{не} влияет на результат, поскольку в кросс-произведениях \eqref{sphere-eq:CS} и в рекурсии \eqref{sphere-eq:recursion} $\chi_n$ входит линейно (нечётно), так что общий множитель сокращается в дробно-линейном отношении. Проверка против аналитического Ми (forward $Q_{\text{sca}}$ и обратное рассеяние $Q_b$) подтверждает корректность знака и того, что замена $h_n^{(1)}\to h_n^{(2)}$ \emph{не} требуется.

\paragraph{Численная устойчивость (PEC-предел).} Для устойчивости при больших $|\operatorname{Im} k_j|$ (толстые поглощающие слои, предел идеального проводника $|\varepsilon|\to\infty$) внутренние функции вычисляются в экспоненциально масштабированной форме \texttt{jve}, \texttt{yve} (масштаб $e^{-|\operatorname{Im} z|}$). Этот масштаб \emph{сокращается} во всех отношениях $\psi'_n/\psi_n$ и кросс-произведениях $C_n,S_n$, так что результат математически тот же, но без переполнения. Производные вычисляются по рекуррентной формуле, следующей из $j_n'(x)=\tfrac{n}{2n+1}j_{n-1}(x)-\tfrac{n+1}{2n+1}j_{n+1}(x)$ \cite{dlmf}:
\begin{equation}
\psi_n'(x)=\frac{n}{2n+1}\,\psi_{n-1}(x)-\frac{n+1}{2n+1}\,\psi_{n+1}(x)+\frac{\psi_n(x)}{x},
\label{sphere-eq:psi-deriv}
\end{equation}
причём общий экспоненциальный масштаб всех трёх членов один и тот же и сокращается при подстановке в отношения (см.\ \texttt{\_psi\_chi}); аналогичная формула справедлива для $\chi_n'$ с заменой $j\to y$.

\subsection{Угловые функции}

Угловая часть выражается через функции $\pi_n,\tau_n$ порядка $m=1$ \cite{bohren1983}:
\begin{align}
\pi_n(\theta) &= \frac{P_n^{1}(\cos\theta)}{\sin\theta}, &
\tau_n(\theta) &= \frac{d P_n^{1}(\cos\theta)}{d\theta}.
\label{sphere-eq:pi-tau}
\end{align}
В коде используются присоединённые функции Лежандра \texttt{scipy.special.lpmv}, включающие фазу Кондона--Шортли $(-1)^m$. С её учётом справедливо тождество \cite{dlmf}
\begin{equation}
\frac{d P_n^{1}(\cos\theta)}{d\theta}=\tfrac12\!\left[n(n+1)\,P_n^{0}(\cos\theta)-P_n^{2}(\cos\theta)\right],
\label{sphere-eq:tau-id}
\end{equation}
тогда как массив \texttt{tay} вычисляется как $\tfrac12\!\left[P_n^{2}-n(n+1)P_n^{0}\right]=-\tau_n$, а массив \texttt{pii} как $P_n^1(\cos\theta)/\sin\theta$ (со сменой знака на $\theta\in(\pi,2\pi)$). Обе угловые функции в коде, таким образом, несут \emph{один и тот же} дополнительный знак $(-1)^1$ относительно конвенции \cite{bohren1983}; поскольку они входят в амплитуды \eqref{sphere-eq:S12} линейно, а диаграмма берётся по модулю, этот общий знак на наблюдаемые величины не влияет. Запись через $P_n^{0},P_n^{1},P_n^{2}$ избегает явного дифференцирования.

\subsection{Послойная рекурсия импеданса (TM) и адмитанса (TE)}

Сферическая симметрия разделяет TM- и TE-каналы, и каждый мультиполь $n$ распространяется \emph{независимо}. Это позволяет ввести для каждой границы поверхностный импеданс $Z_n$ (для TM-мод, ассоциированных с $\bm{N}_{mn}$) и адмитанс $Y_n$ (для TE-мод, ассоциированных с $\bm{M}_{mn}$) и сшивать поля поверхность за поверхностью --- рекурсия Адена--Керкера \cite{aden1951,bohren1983}.

\paragraph{Ядро (слой 0).} В центральном слое регулярно лишь $\psi_n$, поэтому при $k_{00}=k_0 a_0\sigma_0$
\begin{align}
Z_n^{(0)} &= \sqrt{\frac{\varepsilon_1}{\varepsilon_0}}\;\frac{\psi_n'(k_{00})}{\psi_n(k_{00})}, &
Y_n^{(0)} &= \sqrt{\frac{\varepsilon_0}{\varepsilon_1}}\;\frac{\psi_n'(k_{00})}{\psi_n(k_{00})}.
\label{sphere-eq:Z-core}
\end{align}
Корни отношений диэлектрических проницаемостей (для $\mu=1$ это $\sqrt{\varepsilon_{j+1}/\varepsilon_j}$) учитывают скачок волнового импеданса на границе; в общем случае фактор задаётся отношением $\sqrt{\varepsilon\mu}$ для TM и обратным для TE. В коде используется фаза $\varepsilon$ через \texttt{cmath.phase} (atan2), что корректно для $\operatorname{Re}\varepsilon<0$ (металлы, плазмоника).

\paragraph{Шаг рекурсии.} На границе $h$ вводятся кросс-произведения регулярного ($\psi$) и сингулярного ($\chi$) решений по обе стороны границы. Обозначив $A=k_h^{(h)}=k_0 a_h\sigma_h$ (изнутри слоя $h$) и $B=k_{h-1}^{(h)}=k_0 a_{h-1}\sigma_h$ (на предыдущем радиусе в среде слоя $h$):
\begin{align}
C_n &= \psi_n(A)\,\chi_n'(B)-\chi_n(A)\,\psi_n'(B), &
C_n' &= \psi_n'(A)\,\chi_n'(B)-\chi_n'(A)\,\psi_n'(B),\\
S_n &= \chi_n(A)\,\psi_n(B)-\psi_n(A)\,\chi_n(B), &
S_n' &= \chi_n'(A)\,\psi_n(B)-\psi_n'(A)\,\chi_n(B).
\label{sphere-eq:CS}
\end{align}
Тогда импеданс и адмитанс переносятся на следующую границу дробно-линейно \cite{aden1951,bohren1983}:
\begin{align}
Z_n^{(h)} &= \sqrt{\frac{\varepsilon_{h+1}}{\varepsilon_{h}}}\;
\frac{C_n'+Z_n^{(h-1)}\,S_n'}{C_n+Z_n^{(h-1)}\,S_n},
&
Y_n^{(h)} &= \sqrt{\frac{\varepsilon_{h}}{\varepsilon_{h+1}}}\;
\frac{C_n'+Y_n^{(h-1)}\,S_n'}{C_n+Y_n^{(h-1)}\,S_n},
\qquad h=1,\dots,N_L-1.
\label{sphere-eq:recursion}
\end{align}
Рекурсия \eqref{sphere-eq:recursion} --- это перенос логарифмической производной поля через слой; дробно-линейная форма по $Z_n^{(h-1)}$ есть в точности матрица переноса $2\times2$ для пары (поле, его производная), редуцированная к скаляру в силу разделения мод \cite{aden1951,mishchenko2002}. Внешняя среда (воздух, $\varepsilon=1$) дописывается ровно один раз, замыкая рекурсию на внешней границе.

\subsection{Внешние коэффициенты рассеяния и диагональная T-матрица}

На внешней границе $r=R$ (нормировка $R=1$, $z=k_0$) сшивание падающего ($\psi_n$), рассеянного ($\xi_n$) полей и поверхностного импеданса/адмитанса даёт коэффициенты рассеяния \cite{bohren1983,aden1951}:
\begin{align}
M_n &= \frac{Z_n^{(N_L-1)}\,\psi_n(k_0)-\psi_n'(k_0)}{Z_n^{(N_L-1)}\,\xi_n(k_0)-\xi_n'(k_0)},
&
N_n &= \frac{Y_n^{(N_L-1)}\,\psi_n(k_0)-\psi_n'(k_0)}{Y_n^{(N_L-1)}\,\xi_n(k_0)-\xi_n'(k_0)}.
\label{sphere-eq:Mn-Nn}
\end{align}
(В реализации к $M_n,N_n$ применяется комплексное сопряжение --- это переход к конкретной фазовой конвенции диаграммы.) Коэффициенты \eqref{sphere-eq:Mn-Nn} совпадают с классическими коэффициентами Ми $a_n,b_n$ \cite{bohren1983} с точностью до фазовой конвенции, фиксируемой выбором масштаба $\xi_n$ и сопряжением; абсолютная нормировка несущественна, так как все наблюдаемые \eqref{sphere-eq:cross} зависят лишь от $|M_n|^2,|N_n|^2$ и $\operatorname{Re}(M_n+N_n)$ (последнее знакочувствительно, и принятая в коде конвенция выбрана так, чтобы $Q_{\text{ext}}>0$; проверено против аналитического Ми). Существенно лишь то, что для слоистой сферы T-матрица \emph{диагональна} и не смешивает ни порядки $n$, ни TE/TM-каналы, ни азимутальные индексы $m$:
\begin{equation}
\bm{T}=\operatorname{diag}\!\big(M_n,\,N_n\big)_{n\ge1},
\qquad
T_{\sigma l m,\,\sigma' l' m'}=T_{\sigma,l}\,\delta_{\sigma\sigma'}\delta_{ll'}\delta_{mm'}.
\label{sphere-eq:Tmatrix}
\end{equation}
Это и есть формальное выражение того, что многослойная сфера --- задача Ми (Адена--Керкера) с диагональной T-матрицей \cite{aden1951,waterman1971,mishchenko2002}; именно диагональность делает её каноническим ядром для последующей суперпозиции T-матриц (GMM).

\subsection{Поля дальней зоны и сечения}

В дальней зоне рассеянное поле поперечно и определяется амплитудными функциями через ряды по $\pi_n,\tau_n$ \cite{bohren1983}. В терминах коэффициентов \eqref{sphere-eq:Mn-Nn} компоненты диаграммы рассеяния имеют вид
\begin{align}
S_1(\theta)&=\sum_{n\ge1}\frac{2n+1}{n(n+1)}\big(M_n\,\pi_n+N_n\,\tau_n\big),
&
S_2(\theta)&=\sum_{n\ge1}\frac{2n+1}{n(n+1)}\big(M_n\,\tau_n+N_n\,\pi_n\big),
\label{sphere-eq:S12}
\end{align}
откуда для падающей линейно-поляризованной волны
\begin{equation}
E_\theta\propto \cos\varphi\,S_2(\theta),\qquad
E_\varphi\propto -\sin\varphi\,S_1(\theta).
\end{equation}
В коде поле собирается в эквивалентной форме с множителем $(-1)^n$ (фазовая нормировка $(-i)^n$ исходящей волны) и берётся \emph{модуль} комплексной амплитуды (исправление \texttt{real}$\to$\texttt{abs}). Для круговой поляризации со- и кросс-компоненты суть $\propto(\tau_n\mp\pi_n)(M_n\pm N_n)$.

Бистатическое сечение рассеяния (дифференциальное) и интегральные эффективности нормируются на $\pi R^2$ \cite{bohren1983,mishchenko2002}:
\begin{align}
Q_{\text{sca}} &= \frac{2}{x^2}\sum_{n\ge1}(2n+1)\big(|M_n|^2+|N_n|^2\big),\\
Q_{\text{ext}} &= \frac{2}{x^2}\sum_{n\ge1}(2n+1)\,\operatorname{Re}\!\big(M_n+N_n\big),\\
Q_{\text{abs}} &= Q_{\text{ext}}-Q_{\text{sca}},\qquad
Q_{\text{back}} = \frac{1}{x^2}\Big|\sum_{n\ge1}(2n+1)(-1)^n\big(M_n-N_n\big)\Big|^2,
\label{sphere-eq:cross}
\end{align}
где $x=k_0R$ --- параметр размера. Ряды обрезаются по правилу Уискомба $N_{\max}=\lceil x+4x^{1/3}+2\rceil$ \cite{bohren1983,mishchenko2002} (реализовано в \texttt{MieSphere.\_\_post\_init\_\_}). Поглощение учитывается автоматически: комплексный $m=\sqrt{\varepsilon\mu}$ делает $|M_n|^2+|N_n|^2<\operatorname{Re}(M_n+N_n)$ и $Q_{\text{abs}}>0$; предельный переход $|\varepsilon|\to\infty$ (PEC) численно устойчив благодаря масштабированным $\psi_n,\chi_n$.

\subsection{Итог}

Слоистая сфера решается точно для произвольного числа концентрических оболочек: поканальная скалярная рекурсия импеданса/адмитанса \eqref{sphere-eq:Z-core}--\eqref{sphere-eq:recursion} переносит граничные условия от ядра к внешней поверхности, давая коэффициенты \eqref{sphere-eq:Mn-Nn} и диагональную T-матрицу \eqref{sphere-eq:Tmatrix}. Дальнее поле \eqref{sphere-eq:S12} и сечения \eqref{sphere-eq:cross} замыкают задачу. Эта конструкция --- эталон, относительно которого верифицируются сфероидальный и эллипсоидальный решатели и на котором строится T-матричная суперпозиция.

\section{Конфокальный слоистый сфероид}
\label{spheroid-sec:spheroid}

Сфероид (вытянутый --- \emph{prolate}, или сплюснутый --- \emph{oblate} эллипсоид вращения) --- следующий по сложности после сферы конечный объект, для которого скалярное уравнение Гельмгольца разделяется в одной координатной системе \cite{eisenhart1934,flammer1957}. В отличие от сферы, где T-матрица диагональна (см.\ раздел о сфере), сфероидальная задача даёт связь мультипольных порядков внутри фиксированного азимутального индекса $m$. Многослойная задача остаётся строго аналитической \emph{только} при конфокальном вложении оболочек \cite{asano1975,farafonov1996,farafonov2012}. Ниже изложены координаты, разделение, специальные функции, векторные сфероидальные волны, конфокальное условие, рекурсия сшивки по оболочкам и перепроекция в сферический базис VSWF для последующего объединения по схеме GMM.

\subsection{Сфероидальные координаты и параметр размера}

Введём фокусное расстояние $d$ (расстояние между фокусами для prolate; диаметр фокального кольца для oblate) и безразмерный параметр размера
\begin{equation}
c \;=\; \frac{k d}{2}, \qquad k = \omega\sqrt{\varepsilon\mu},
\label{spheroid-eq:cdef}
\end{equation}
где $k$ --- волновое число в данной среде (в общем случае комплексное для поглощающего слоя). Для \emph{вытянутого} сфероида декартовы координаты выражаются через $(\xi,\eta,\varphi)$, $\xi\in[1,\infty)$, $\eta\in[-1,1]$, $\varphi\in[0,2\pi)$:
\begin{align}
x &= \tfrac{d}{2}\sqrt{(\xi^2-1)(1-\eta^2)}\cos\varphi, \nonumber\\
y &= \tfrac{d}{2}\sqrt{(\xi^2-1)(1-\eta^2)}\sin\varphi, \label{spheroid-eq:prolate-cart}\\
z &= \tfrac{d}{2}\,\xi\eta. \nonumber
\end{align}
Поверхность $\xi=\xi_0=\mathrm{const}$ есть вытянутый сфероид с полуосями $a=\tfrac{d}{2}\xi_0$ (вдоль оси $z$) и $b=\tfrac{d}{2}\sqrt{\xi_0^2-1}$ (поперёк), так что
\begin{equation}
a^2 - b^2 = (d/2)^2 \qquad(\text{prolate}).
\label{spheroid-eq:prolate-focal}
\end{equation}
Для \emph{сплюснутого} сфероида делается замена $\xi\to i\xi$ (эквивалентно $c\to -ic$, $\xi\in[0,\infty)$) \cite{flammer1957,dlmf}:
\begin{align}
x &= \tfrac{d}{2}\sqrt{(\xi^2+1)(1-\eta^2)}\cos\varphi, \nonumber\\
y &= \tfrac{d}{2}\sqrt{(\xi^2+1)(1-\eta^2)}\sin\varphi, \label{spheroid-eq:oblate-cart}\\
z &= \tfrac{d}{2}\,\xi\eta, \nonumber
\end{align}
с $a=\tfrac{d}{2}\sqrt{\xi_0^2+1}$, $b=\tfrac{d}{2}\xi_0$ (теперь $a$ --- экваториальная, $b$ --- полярная полуось) и
\begin{equation}
a^2 - b^2 = (d/2)^2 \qquad(\text{oblate}).
\label{spheroid-eq:oblate-focal}
\end{equation}
Метрические коэффициенты Ламе (prolate) равны
\begin{equation}
h_\xi = \frac{d}{2}\sqrt{\frac{\xi^2-\eta^2}{\xi^2-1}},\quad
h_\eta = \frac{d}{2}\sqrt{\frac{\xi^2-\eta^2}{1-\eta^2}},\quad
h_\varphi = \frac{d}{2}\sqrt{(\xi^2-1)(1-\eta^2)},
\label{spheroid-eq:lame}
\end{equation}
а элемент объёма $dV = h_\xi h_\eta h_\varphi\,d\xi\,d\eta\,d\varphi$. Предел $d\to0$ при фиксированном $a$ ($\xi_0\to\infty$, $\tfrac{d}{2}\xi_0\to a$) восстанавливает сферические координаты $(r,\theta,\varphi)$, $r=\tfrac{d}{2}\xi$, $\eta=\cos\theta$.

\subsection{Разделение скалярного уравнения Гельмгольца}

Скалярное уравнение $(\nabla^2 + k^2)\Psi = 0$ при подстановке $\Psi = R_{ml}(c,\xi)\,S_{ml}(c,\eta)\,e^{im\varphi}$ распадается на угловое и радиальное сфероидальные уравнения \cite{flammer1957,stratton1941,dlmf}:
\begin{align}
\frac{d}{d\eta}\!\left[(1-\eta^2)\frac{dS_{ml}}{d\eta}\right]
&+\left(\lambda_{ml} - c^2\eta^2 - \frac{m^2}{1-\eta^2}\right)S_{ml}=0,
\label{spheroid-eq:angular}\\[4pt]
\frac{d}{d\xi}\!\left[(\xi^2-1)\frac{dR_{ml}}{d\xi}\right]
&-\left(\lambda_{ml} - c^2\xi^2 + \frac{m^2}{\xi^2-1}\right)R_{ml}=0,
\label{spheroid-eq:radial}
\end{align}
где $\lambda_{ml}(c)$ --- общая постоянная разделения (собственное значение), а $l\ge m$ нумерует решения углового уравнения. Для oblate заменяют $\xi^2-1\to\xi^2+1$ и $c^2\to -c^2$ в \eqref{spheroid-eq:radial} \cite{flammer1957,dlmf}.

\paragraph{Угловые функции.} Решения \eqref{spheroid-eq:angular}, регулярные на $\eta=\pm1$, --- угловые сфероидальные функции $S_{ml}(c,\eta)$, разлагаемые по присоединённым полиномам Лежандра:
\begin{equation}
S_{ml}(c,\eta)=\sum_{r=0,1}^{\infty}{}' \, d_r^{ml}(c)\,P_{m+r}^{m}(\eta),
\label{spheroid-eq:Sexpand}
\end{equation}
где штрих означает суммирование по $r$ той же чётности, что $l-m$, а коэффициенты $d_r^{ml}(c)$ удовлетворяют трёхчленной рекурсии, определяющей одновременно $\lambda_{ml}(c)$ как собственное значение \cite{flammer1957,dlmf}. При $c\to0$ имеем $\lambda_{ml}\to l(l+1)$ и $S_{ml}(0,\eta)\to P_l^m(\eta)$.

\paragraph{Радиальные функции.} Решения \eqref{spheroid-eq:radial} --- радиальные сфероидальные функции $R_{ml}^{(j)}(c,\xi)$, $j=1,2,3,4$, представимые рядами по сферическим функциям Бесселя $z_n^{(j)}$ от аргумента $c\xi$ \cite{flammer1957,dlmf}:
\begin{equation}
R_{ml}^{(j)}(c,\xi)=
\frac{\left(\frac{\xi^2-1}{\xi^2}\right)^{m/2}}{\sum'_r d_r^{ml}\frac{(2m+r)!}{r!}}
\sum_{r=0,1}^{\infty}{}'\, i^{\,r+m-l}\,\frac{(2m+r)!}{r!}\,d_r^{ml}(c)\,z_{m+r}^{(j)}(c\xi),
\label{spheroid-eq:Rexpand}
\end{equation}
с $z_n^{(1)}=j_n$ (регулярна в начале координат), $z_n^{(2)}=y_n$ и
\begin{equation}
R_{ml}^{(3)} = R_{ml}^{(1)} + i\,R_{ml}^{(2)}\;\;\xrightarrow[\xi\to\infty]{}\;\;
\frac{e^{i(c\xi - (l+1)\pi/2)}}{c\xi},
\label{spheroid-eq:R3}
\end{equation}
что в принятой конвенции $e^{-i\omega t}$ задаёт \emph{исходящую} волну (полный аналог $h_n^{(1)}$ сферической задачи; ср.\ $z_n^{(3)}=h_n^{(1)}$). $R_{ml}^{(1)}$ описывает регулярное (внутреннее/падающее) поле. Таким образом разделение в $(\xi,\eta,\varphi)$ полностью параллельно сферическому $\{j_n, h_n^{(1)}, Y_l^m\}$, с заменой $j_n\!\to\!R^{(1)}_{ml}$, $h_n^{(1)}\!\to\!R^{(3)}_{ml}$, $P_l^m\!\to\!S_{ml}$.

\subsection{Векторные сфероидальные волновые функции}

Векторные сфероидальные волновые функции (VSWF) строятся как у Стрэттона \cite{stratton1941,asano1975}. Из порождающих скаляров $\psi_{ml}^{(j)}=R_{ml}^{(j)}(c,\xi)\,S_{ml}(c,\eta)\,e^{im\varphi}$ образуются
\begin{equation}
\bm{M}_{ml}^{(j)}=\nabla\times(\bm{r}\,\psi_{ml}^{(j)}),\qquad
\bm{N}_{ml}^{(j)}=\frac{1}{k}\,\nabla\times\bm{M}_{ml}^{(j)},
\label{spheroid-eq:MN}
\end{equation}
где $\bm{M}$ соответствует TE (магнитному, поперечно-электрическому относительно $\bm r$), $\bm{N}$ --- TM (электрическому) типу. Полное поле в слое раскладывается по $\{\bm{M}_{ml}^{(1)},\bm{N}_{ml}^{(1)},\bm{M}_{ml}^{(3)},\bm{N}_{ml}^{(3)}\}$. В отличие от сферы, где $\bm M,\bm N$ диагональны по $l$, здесь координатные орты $(\hat{\bm\xi},\hat{\bm\eta},\hat{\bm\varphi})$ перемешивают компоненты: при фиксированном $m$ (азимутальная зависимость $e^{im\varphi}$ сохраняется и развязывает разные $m$) разные мультипольные порядки $l$ оказываются \emph{связанными} как через геометрию орт, так и через сшивку TE/TM. Развязка $\bm M$ и $\bm N$ (чистые TE/TM) имеет место только при осевом падении ($m=0$, волна вдоль оси симметрии) \cite{asano1975,farafonov2012}.

\subsection{Конфокальное многослойное условие}

Пусть оболочки занумерованы $j=1,\dots,L$ (внутреннее ядро $j=1$), с границами $\xi=\xi_j$. Аналитическое разделение во всех слоях одновременно возможно \emph{тогда и только тогда}, когда все границы --- софокусные сфероиды с \emph{общим} фокусным расстоянием $d$:
\begin{equation}
a_j^2 - b_j^2 = (d/2)^2 \qquad \forall\, j=1,\dots,L,
\label{spheroid-eq:confocal}
\end{equation}
--- прямой аналог условия Фарафонова--Вощинникова \cite{farafonov2012,farafonov1996}. Из \eqref{spheroid-eq:confocal} следует, что геометрический параметр $\xi_j$ \emph{одинаков} для системы, а форм-фактор (аспектное отношение $a_j/b_j$) меняется от слоя к слою: внутренние оболочки более вытянуты/сплюснуты. При этом параметр размера зависит от слоя через локальное волновое число:
\begin{equation}
c_j = \tfrac12 k_j d = \tfrac12 \omega \sqrt{\varepsilon_j \mu_j}\, d,
\label{spheroid-eq:cj}
\end{equation}
так что базисы соседних слоёв построены на \emph{разных} $c_j$ --- это ключевое отличие от сферы, порождающее интегралы перекрытия (см.\ ниже). Внешняя среда (вакуум/воздух) имеет $c_0 = \tfrac12 k_0 d$.

\subsection{Сшивка на границах и блочно-матричная рекурсия}

На каждой границе $\xi=\xi_j$ выполняются условия непрерывности тангенциальных $\bm E$ и $\bm H$. Поле в слое $j$ записывается через коэффициенты разложения; для фиксированного $m$ соберём амплитуды по $l$ в векторы-столбцы $\bm{a}_j^{M},\bm{a}_j^{N}$ (регулярная часть, $R^{(1)}$) и $\bm{b}_j^{M},\bm{b}_j^{N}$ (исходящая, $R^{(3)}$). Сшивка даёт линейную систему, в которой индекс $m$ \emph{развязан} (отдельная задача на каждое $m$), а порядки $l$ \emph{связаны} --- т.е.\ работаем с блочными матрицами размера $2N_l\times2N_l$ (множитель $2$ --- сцепление TE/TM, $N_l$ --- число удерживаемых $l$):
\begin{equation}
\begin{pmatrix}\bm{a}_{j+1}\\ \bm{b}_{j+1}\end{pmatrix}
=
\bm{T}_j^{(m)}
\begin{pmatrix}\bm{a}_{j}\\ \bm{b}_{j}\end{pmatrix},
\qquad
\bm{T}_j^{(m)} =
\big[\bm{Q}_{j+1}^{(m)}(\xi_j)\big]^{-1}\,\bm{Q}_{j}^{(m)}(\xi_j),
\label{spheroid-eq:transfer}
\end{equation}
где $\bm{Q}_{j}^{(m)}(\xi_j)$ --- блочная матрица граничных функций слоя $j$, составленная из $R_{ml}^{(1,3)}(c_j,\xi_j)$, их $\xi$-производных и связанных с $\bm M,\bm N$ комбинаций (тангенциальные проекции на $\xi=\xi_j$). Эта матрица --- прямое обобщение скалярных кросс-произведений $C,S$ и импедансов $Z,Y$ сферической рекурсии Адена--Керкера: в сферическом пределе $\bm{Q}^{(m)}$ становится блочно-диагональной по $l$ и распадается на независимые $2\times2$ блоки (TE/TM), воспроизводя $Z_n=\sqrt{\varepsilon_{j+1}/\varepsilon_j}\,(\cdots)$ из ядра \texttt{mie\_core}.

\paragraph{Внутреннее замыкание и внешний коэффициент.} В ядре $j=1$ исходящие амплитуды отсутствуют ($\bm{b}_1=\bm0$, регулярность в начале), что задаёт начальное условие. Перемножение трансфер-блоков
\begin{equation}
\bm{T}^{(m)} = \bm{T}_{L}^{(m)}\,\bm{T}_{L-1}^{(m)}\cdots\bm{T}_{1}^{(m)}
\label{spheroid-eq:product}
\end{equation}
даёт связь поля ядра с полем во внешней среде; размер $\bm{T}^{(m)}$ \emph{не зависит} от числа слоёв $L$ (накопление в произведение), как и в скалярном случае. Согласование с разложением падающей плоской волны по $\{\bm M^{(1)},\bm N^{(1)}\}$ внешней среды и требование, чтобы рассеянное поле содержало только $\{\bm M^{(3)},\bm N^{(3)}\}$, замыкает систему и даёт сфероидальную T-матрицу
\begin{equation}
\begin{pmatrix}\bm{p}^{M}\\ \bm{p}^{N}\end{pmatrix}_{\!\!\text{sca}}
=
\bm{T}_{\text{sph}}^{(m)}
\begin{pmatrix}\bm{a}^{M}\\ \bm{a}^{N}\end{pmatrix}_{\!\!\text{inc}},
\label{spheroid-eq:Tspheroid}
\end{equation}
недиагональную по $l$ (блочную) для каждого $m$ \cite{asano1975,farafonov2012}.

\paragraph{Интегралы перекрытия между слоями.} Поскольку $c_j\ne c_{j+1}$ (разные $k_j$), угловые функции соседних слоёв образуют \emph{разные} неортогональные базисы. Сшивка тангенциальных полей на $\xi=\xi_j$ требует разложения $S_{ml}(c_j,\eta)$ по $\{S_{ml'}(c_{j+1},\eta)\}$, что даёт матрицу перекрытия
\begin{equation}
G^{(m)}_{ll'}(c_j,c_{j+1})=\frac{\displaystyle\int_{-1}^{1} S_{ml}(c_j,\eta)\,S_{ml'}(c_{j+1},\eta)\,d\eta}
{\displaystyle\int_{-1}^{1} S_{ml'}^2(c_{j+1},\eta)\,d\eta},
\label{spheroid-eq:overlap}
\end{equation}
которая входит множителем в граничные блоки $\bm{Q}_j^{(m)}$ и связывает порядки $l$. На практике $G^{(m)}_{ll'}$ вычисляется через коэффициенты $d_r^{ml}(c)$ из \eqref{spheroid-eq:Sexpand} и интегралы произведений присоединённых функций Лежандра \cite{farafonov2012,flammer1957}. При $c_j=c_{j+1}$ (однородная среда) $G^{(m)}_{ll'}=\delta_{ll'}$, и связь $l$ исчезает.

\subsection{Перепроекция в сферический базис VSWF для GMM}

Для объединения сфероида с другими примитивами по схеме суперпозиции T-матриц (GMM) сфероидальная T-матрица \eqref{spheroid-eq:Tspheroid} переводится в \emph{стандартный} сферический базис VSWF $\{\bm M_{nm},\bm N_{nm}\}$. Используется разложение сфероидальных VSWF по сферическим с матрицами связи $\alpha^{(m)}_{nl}(c)$, $\beta^{(m)}_{nl}(c)$ (получаемыми из \eqref{spheroid-eq:Sexpand}--\eqref{spheroid-eq:Rexpand}, поскольку $R^{(j)}_{ml}$ уже выражены через $z_n^{(j)}(c\xi)$, а $S_{ml}$ --- через $P_n^m$):
\begin{equation}
\bm{M}_{ml}^{(j)}(c)=\sum_{n\ge m}\Big[\alpha^{(m)}_{nl}\,\bm{M}_{nm}^{(j)} + \beta^{(m)}_{nl}\,\bm{N}_{nm}^{(j)}\Big],
\label{spheroid-eq:reproject}
\end{equation}
и аналогично для $\bm N_{ml}^{(j)}$. Собрав $\alpha,\beta$ в блочную матрицу $\bm{P}^{(m)}(c_0)$ (на внешнем параметре $c_0$), получаем сферическую T-матрицу
\begin{equation}
\bm{T}^{(m)}_{\text{sphere-basis}}=\bm{P}^{(m)}(c_0)\,\bm{T}_{\text{sph}}^{(m)}\,\big[\bm{P}^{(m)}(c_0)\big]^{-1},
\label{spheroid-eq:Tsphere}
\end{equation}
недиагональную по $n$ (в отличие от диагональной Mie-матрицы сферы) --- именно её передают в аддитивные теоремы Стейна--Крузана для трансляций между центрами рассеивателей в GMM \cite{stein1961,cruzan1962,xu1995,waterman1971}.

\subsection{Дальняя зона и сечения}

Из исходящей части $R^{(3)}_{ml}$ и асимптотики \eqref{spheroid-eq:R3} получается амплитудная матрица рассеяния и, через оптическую теорему, сечения. В стандартной T-матричной нормировке (конвенция $e^{-i\omega t}$, $\bm{p}=\bm{T}\bm{a}$) ориентационно-усреднённые сечения выражаются через след T-матрицы \cite{mishchenko2002,waterman1971}:
\begin{equation}
C_{\text{ext}}=-\frac{2\pi}{k_0^2}\,\mathrm{Re}\!\sum_{m}\mathrm{Tr}\big[\bm{T}^{(m)}_{\text{sphere-basis}}\big],\qquad
C_{\text{sca}}=\frac{2\pi}{k_0^2}\sum_{m}\mathrm{Tr}\big[\bm{T}^{(m)\dagger}_{\text{sphere-basis}}\,\bm{T}^{(m)}_{\text{sphere-basis}}\big],
\label{spheroid-eq:cross}
\end{equation}
где $\mathrm{Tr}[\bm T^{\dagger}\bm T]=\sum_{nn'}|T_{nn'}|^2$ --- квадрат фробениусовой нормы блока. Знак минус в $C_{\text{ext}}$ и множитель $2\pi/k_0^2$ --- стандартная нормировка оптической теоремы для T-матрицы \cite{mishchenko2002}; для фиксированной ориентации $C_{\text{ext}}=-(4\pi/k_0^2)\,\mathrm{Re}[\bm{a}^{(m)\dagger}_{\text{inc}}\,\bm{T}^{(m)}\bm{a}^{(m)}_{\text{inc}}]$, $C_{\text{sca}}=(4\pi/k_0^2)\,\|\bm{T}^{(m)}\bm{a}^{(m)}_{\text{inc}}\|^2$ при единично нормированных амплитудах падающей волны \cite{mishchenko2002,asano1975}. Поглощение --- $C_{\text{abs}}=C_{\text{ext}}-C_{\text{sca}}>0$ для поглощающих слоёв.

\subsection{Предельный переход к сфере ($d\to0$)}

При $d\to0$ ($c\to0$) угловые функции $S_{ml}(0,\eta)=P_l^m(\eta)$, радиальные $R^{(1)}_{ml}(0,\xi)\propto j_l(kr)$, $R^{(3)}_{ml}(0,\xi)\propto h_l^{(1)}(kr)$ \cite{flammer1957,dlmf}. Матрица перекрытия $G^{(m)}_{ll'}\to\delta_{ll'}$, блок $\bm{Q}^{(m)}$ становится блочно-диагональным, связь $l$ исчезает, TE/TM полностью разделяются, и блочная рекурсия \eqref{spheroid-eq:transfer} вырождается в скалярную импеданс/адмитанс-рекурсию слоистой сферы (Аден--Керкер) \cite{aden1951,bohren1983} --- то самое ядро \texttt{mie\_core} с $Z_n,Y_n$ и кросс-произведениями $C_n,S_n$. T-матрица \eqref{spheroid-eq:Tsphere} становится диагональной (Mie). Это даёт прямой численный контроль реализации.

\subsection{Численный аспект: комплексный $c$ для поглощающих слоёв}

Для поглощающего слоя $\varepsilon_j$ комплексна, поэтому $c_j=\tfrac12\omega\sqrt{\varepsilon_j\mu_j}\,d$ --- комплексный параметр размера. Вычисление $\lambda_{ml}(c)$, коэффициентов $d_r^{ml}(c)$ и радиальных функций $R^{(j)}_{ml}(c,\xi)$ при комплексном $c$ --- известная численно тонкая задача: трёхчленная рекурсия для $d_r^{ml}$ теряет устойчивость, а ряды \eqref{spheroid-eq:Rexpand} страдают от катастрофического сокращения, особенно для oblate с большим $|\mathrm{Im}\,c|$ \cite{vanburen2020}. Требуются устойчивые алгоритмы (непрерывные дроби для $\lambda_{ml}$, контролируемое прямое/обратное направление рекурсии, расширенная точность), как в работе Ван Бюрена и Буавера \cite{vanburen2020}; это прямой сфероидальный аналог экспоненциально масштабированных функций Риккати--Бесселя (\texttt{jve}/\texttt{yve}), применённых в сферическом ядре для устойчивости при больших $|\mathrm{Im}\,(k a)|$.

\section{Триаксиальный эллипсоид}
\label{ellipsoid-sec:ellipsoid}

Триаксиальный (трёхосный) эллипсоид с полуосями $a>b>c$,
\begin{equation}
\frac{x^2}{a^2}+\frac{y^2}{b^2}+\frac{z^2}{c^2}=1,
\label{ellipsoid-eq:ell-surface}
\end{equation}
замыкает иерархию конечных тел, ограниченных одной координатной поверхностью: сфера и сфероид суть его вырожденные случаи ($a=b=c$ и $a=b\neq c$). Это последняя из систем, в которых скалярное уравнение Гельмгольца разделяется в смысле Эйзенхарта \cite{eisenhart1934}. Однако, как показано ниже, переход от скалярной задачи к \emph{полной векторной} задаче Максвелла здесь принципиально иной, чем для сферы и сфероида: практического разделяющегося решения вектора не существует. Поэтому раздел построен в два этапа --- сначала фиксируется граница применимости (скалярное разделение и его векторная непригодность), затем даётся реально работающий аналитический результат --- квазистатическая (рэлеевская) многослойная конфокальная задача \cite{sihvola1990,bohren1983}.

\subsection{Эллипсоидальные координаты и конфокальные квадрики}

Введём эллипсоидальные координаты $(\xi,\eta,\zeta)$ как корни кубического (по параметру $u$) уравнения семейства конфокальных квадрик с фокальными параметрами $h^2=a^2-c^2$, $k^2=a^2-b^2$ (так что $b^2-c^2=h^2-k^2$) \cite{stratton1941}:
\begin{equation}
\frac{x^2}{u}+\frac{y^2}{u-k^2}+\frac{z^2}{u-h^2}=1 .
\label{ellipsoid-eq:confocal-family}
\end{equation}
Три корня попадают в непересекающиеся интервалы и образуют ортогональную систему координат:
\begin{equation}
\xi\in[h^2,\infty)\ (\text{эллипсоиды}),\quad
\eta\in[k^2,h^2]\ (\text{однополостные гиперболоиды}),\quad
\zeta\in[0,k^2]\ (\text{двуполостные гиперболоиды}).
\label{ellipsoid-eq:coord-ranges}
\end{equation}
Поверхность частицы~\eqref{ellipsoid-eq:ell-surface} есть координатная поверхность $\xi=a^2$ (при $u=a^2$ в~\eqref{ellipsoid-eq:confocal-family} получаем именно~\eqref{ellipsoid-eq:ell-surface}). Декартовы координаты выражаются через $(\xi,\eta,\zeta)$ рационально-радикально, например
\begin{equation}
x^2=\frac{\xi\,\eta\,\zeta}{k^2 h^2},\qquad
y^2=\frac{(\xi-k^2)(\eta-k^2)(\zeta-k^2)}{k^2(k^2-h^2)},\qquad
z^2=\frac{(\xi-h^2)(\eta-h^2)(\zeta-h^2)}{h^2(h^2-k^2)} .
\label{ellipsoid-eq:cart-from-ell}
\end{equation}
Метрические коэффициенты Ламе для~\eqref{ellipsoid-eq:confocal-family} имеют вид
\begin{equation}
h_\xi^2=\frac{(\xi-\eta)(\xi-\zeta)}{4\,R(\xi)},\quad
R(\xi)=\xi(\xi-k^2)(\xi-h^2),
\label{ellipsoid-eq:metric-xi}
\end{equation}
и циклически для $\eta,\zeta$ (с тем же кубическим радикандом $R(\cdot)$). Существенно, что одна координатная поверхность $\xi=\mathrm{const}$ исчерпывает геометрию: именно это делает эллипсоид допустимым «примитивом» наряду со сферой и сфероидом.

\subsection{Разделение скалярного уравнения: функции Ламе}

Подстановка $\Psi(\xi,\eta,\zeta)=E(\xi)\,E(\eta)\,E(\zeta)$ в уравнение Гельмгольца $(\nabla^2+\kappa^2)\Psi=0$ с метрикой~\eqref{ellipsoid-eq:metric-xi} приводит к разделению на три формально \emph{идентичных} обыкновенных дифференциальных уравнения --- уравнения Ламе в алгебраической (форме Якоби) форме \cite{stratton1941}:
\begin{equation}
4\,R(u)\,\frac{d^2 E}{du^2}+2\,R'(u)\,\frac{dE}{du}-\bigl(\kappa^2 u^2+\alpha_1 u+\alpha_0\bigr)E(u)=0,
\qquad u\in\{\xi,\eta,\zeta\},
\label{ellipsoid-eq:lame-wave}
\end{equation}
где $R(u)=u(u-k^2)(u-h^2)$, а $\alpha_0,\alpha_1$ --- две константы разделения, общие для всех трёх множителей. В пределе $\kappa\to0$ (Лаплас, статика) член $\kappa^2 u^2$ исчезает, и~\eqref{ellipsoid-eq:lame-wave} переходит в \emph{классическое} уравнение Ламе, регулярные полиномиальные решения которого --- функции Ламе (эллипсоидальные гармоники) $E_n^m(u)$ степени $n$, образующие при $1\le m\le 2n+1$ полный набор из $2n+1$ функций. Произведения $E_n^m(\xi)E_n^m(\eta)E_n^m(\zeta)$ суть поверхностные эллипсоидальные гармоники; внешние (убывающие) решения строятся через сопряжённые функции Ламе второго рода
\begin{equation}
F_n^m(\xi)=(2n+1)\,E_n^m(\xi)\!\int_\xi^\infty\frac{du}{\bigl[E_n^m(u)\bigr]^2\sqrt{R(u)}} .
\label{ellipsoid-eq:lame-second-kind}
\end{equation}
При $\kappa\neq0$ полные решения~\eqref{ellipsoid-eq:lame-wave} --- \emph{волновые} функции Ламе --- уже не полиномиальны и не сводятся к замкнутым специальным функциям; их приходится строить численно как ряды по сфероидальным/сферическим базисам. Это и есть техническая причина следующего пункта.

\subsection{Граница применимости: векторная задача не разделяется}

Для сферы и сфероида скалярное разделение поднимается до векторного: применяя к скалярным решениям пилотный вектор и оператор $\nabla\times$, строят соленоидальные векторные сферические/сфероидальные волновые функции $\bm{M},\bm{N}$, на которых граничные условия Максвелла замыкаются в (диагональную для сферы, блочную по $m$ для сфероида) систему. Для триаксиального эллипсоида эта конструкция \emph{обрывается}: оператор Гельмгольца разделяется, но векторное уравнение $\nabla\times\nabla\times\bm{E}-\kappa^2\bm{E}=0$ не допускает построения полного семейства поперечных решений, разделяющихся в координатах~\eqref{ellipsoid-eq:confocal-family}. Эллипсоидальный пилотный вектор не порождает решений, удовлетворяющих $\nabla\!\cdot\bm{E}=0$ в той же сепарабельной форме, и три уравнения Ламе~\eqref{ellipsoid-eq:lame-wave} не расцепляют поляризационные компоненты. Дополнительно сами волновые функции Ламе вычислительно неподъёмны вне малых степеней $n$, поэтому зрелого, верифицированного полноволнового сепарабельного решателя для произвольного $\kappa a$ не существует \cite{stratton1941,bohren1983}. Таким образом, в нашем семействе примитивов триаксиальный эллипсоид --- единственное тело, для которого полноволновой аналитический модуль \emph{не} строится; общий путь к произвольному размеру для него --- численные методы (T-матрица/EBCM, объёмные интегральные уравнения), выходящие за рамки данной аналитической библиотеки. Аналитически замкнутый многослойный результат для эллипсоида существует только в квазистатике, к которой переходим.

\subsection{Квазистатика (Рэлей): деполяризационные коэффициенты}

В пределе малого параметра размера, $x=\kappa\,a\ll1$ (и $|m|\kappa a\ll1$ внутри), запаздыванием можно пренебречь: внутри и вблизи частицы поле потенциально, $\bm{E}=-\nabla\phi$, $\nabla^2\phi=0$ (Лаплас). Это снимает все трудности предыдущего пункта --- уравнение~\eqref{ellipsoid-eq:lame-wave} вырождается в классическое уравнение Ламе, и для однородного эллипсоида в однородном внешнем поле $\bm{E}_0$ достаточно гармоник степени $n=1$. Решение даёт постоянное внутреннее поле $\bm{E}_{\mathrm{in}}$, связанное с приложенным покомпонентно вдоль главных осей $j\in\{a,b,c\}$ \cite{bohren1983,stratton1941}:
\begin{equation}
E_{\mathrm{in},j}=\frac{E_{0,j}}{1+L_j\,(\varepsilon_p/\varepsilon_m-1)},
\label{ellipsoid-eq:internal-field}
\end{equation}
где $\varepsilon_p$ --- проницаемость частицы, $\varepsilon_m$ --- среды, а $L_j$ --- \emph{геометрические деполяризационные коэффициенты}. Они задаются эллиптическими интегралами по внешней эллипсоидальной координате
\begin{equation}
L_j=\frac{abc}{2}\int_0^\infty\frac{ds}{(s+a_j^2)\,\sqrt{(s+a^2)(s+b^2)(s+c^2)}},
\qquad (a_a,a_b,a_c)\equiv(a,b,c),
\label{ellipsoid-eq:depol-integral}
\end{equation}
и удовлетворяют тождеству суммы (след тензора деполяризации единичен)
\begin{equation}
L_a+L_b+L_c=1,\qquad 0<L_j<1 .
\label{ellipsoid-eq:depol-sum}
\end{equation}
Интеграл~\eqref{ellipsoid-eq:depol-integral} приводится к стандартным неполным эллиптическим интегралам Лежандра $F(\varphi,k)$, $E(\varphi,k)$ с амплитудой и модулем
\begin{equation}
\varphi=\arcsin\!\sqrt{1-\frac{c^2}{a^2}}=\arccos\frac{c}{a},\qquad
k=\sqrt{\frac{a^2-b^2}{a^2-c^2}} ,
\label{ellipsoid-eq:elliptic-params}
\end{equation}
например для наибольшей оси
\begin{equation}
L_a=\frac{abc}{(a^2-b^2)\sqrt{a^2-c^2}}\bigl[F(\varphi,k)-E(\varphi,k)\bigr].
\label{ellipsoid-eq:La-explicit}
\end{equation}
Предельные случаи фиксируют согласованность: сфера $a=b=c\Rightarrow L_a=L_b=L_c=1/3$; вытянутый сфероид $a>b=c\Rightarrow L_b=L_c=(1-L_a)/2$; иголка $L_a\to0$, $L_b=L_c\to1/2$; диск $L_a\to1$, $L_b=L_c\to0$ \cite{bohren1983}.

\subsection{Поляризуемость однородного эллипсоида}

Индуцированный дипольный момент вдоль оси $j$ есть $p_j=\varepsilon_0\varepsilon_m\,\alpha_j\,E_{0,j}$ с главными поляризуемостями (тензор $\bm{\alpha}$ диагонален в осях эллипсоида) \cite{bohren1983,stratton1941}:
\begin{equation}
\alpha_j=V\,\frac{\varepsilon_p-\varepsilon_m}{\varepsilon_m+L_j\,(\varepsilon_p-\varepsilon_m)},\qquad
V=\frac{4}{3}\pi abc .
\label{ellipsoid-eq:alpha-homog}
\end{equation}
При $L_j=1/3$ это в точности рэлеевский предел сферы $\alpha=3V\,(\varepsilon_p-\varepsilon_m)/(\varepsilon_p+2\varepsilon_m)$. Полюс знаменателя при $\varepsilon_p=-\varepsilon_m(1-L_j)/L_j$ --- условие фрёлиховского (плазмонного) резонанса вдоль оси $j$; три различных $L_j$ дают три раздельных резонанса --- характерную «расщеплённую» спектральную сигнатуру триаксиальной частицы.

\subsection{Конфокальная многослойная частица: матрицы переноса Сиволы--Линделла}

Аналитическое продолжение на \emph{многослойный} эллипсоид требует, чтобы все границы слоёв были \emph{конфокальны}, т.е. являлись поверхностями $\xi=\mathrm{const}$ одного и того же семейства~\eqref{ellipsoid-eq:confocal-family}: фокальные параметры $h^2,k^2$ общие для всех оболочек, а полуоси $p$-го слоя удовлетворяют
\begin{equation}
a_p^2-b_p^2=k^2=\mathrm{const},\qquad a_p^2-c_p^2=h^2=\mathrm{const}\quad(\forall p),
\label{ellipsoid-eq:confocal-req}
\end{equation}
что эквивалентно $a_p^2=a_q^2+\delta_{pq}$, $b_p^2=b_q^2+\delta_{pq}$, $c_p^2=c_q^2+\delta_{pq}$ с общим сдвигом $\delta_{pq}$ для всех трёх осей. Только при~\eqref{ellipsoid-eq:confocal-req} граничные условия Лапласа на каждой границе разделяются в тех же координатах, и $n=1$-гармоника не перемешивается с высшими; форма (соотношение осей) при этом необходимо меняется от слоя к слою. Для каждой главной оси $j$ потенциал в слое $p$ есть суперпозиция «растущей» ($\propto E_1$) и «убывающей» ($\propto F_1$, см.~\eqref{ellipsoid-eq:lame-second-kind}) гармоник степени $1$; их амплитуды $(A_p,B_p)^{\mathsf T}$ связываются на границе $\xi_p$ (непрерывность потенциала и нормальной компоненты $\varepsilon\,\partial\phi/\partial n$) линейно, что и даёт $2\times2$ \emph{матрицу переноса} Сиволы--Линделла \cite{sihvola1990}.

Для двухслойной частицы (ядро $\varepsilon_1$ с деполяризационными коэффициентами $L_j^{(1)}$ его границы; внешняя оболочка $\varepsilon_2$ с коэффициентами $L_j^{(2)}$ её внешней границы; объёмная доля ядра $f=V_1/V_2$, $V_2=\tfrac43\pi a_2 b_2 c_2$) последовательная сшивка потенциала даёт замкнутую поляризуемость в среде $\varepsilon_m$ вдоль оси $j$ \cite{sihvola1990,bohren1983}:
\begin{equation}
\alpha_j^{(2)}=V_2\,
\frac{\bigl(\varepsilon_2-\varepsilon_m\bigr)\bigl[\varepsilon_2+(\varepsilon_1-\varepsilon_2)\bigl(L_j^{(1)}-f L_j^{(2)}\bigr)\bigr]
      +f\,\varepsilon_2\,(\varepsilon_1-\varepsilon_2)}
     {\bigl[\varepsilon_2+(\varepsilon_1-\varepsilon_2)\bigl(L_j^{(1)}-f L_j^{(2)}\bigr)\bigr]
      \bigl[\varepsilon_m+(\varepsilon_2-\varepsilon_m)L_j^{(2)}\bigr]
      +f\,L_j^{(2)}\,\varepsilon_2\,(\varepsilon_1-\varepsilon_2)} ,
\label{ellipsoid-eq:two-layer-alpha}
\end{equation}
где $L_j^{(1)}$ --- коэффициент внутренней (ядерной) границы, $L_j^{(2)}$ --- внешней. При $L_j^{(1)}=L_j^{(2)}=1/3$ формула~\eqref{ellipsoid-eq:two-layer-alpha} тождественно сводится к классической поляризуемости двухслойной (концентрической) сферы Адена--Керкера \cite{aden1951,bohren1983}, а при $f\to0$ или $\varepsilon_1=\varepsilon_2$ --- к однослойной~\eqref{ellipsoid-eq:alpha-homog}. Эквивалентно~\eqref{ellipsoid-eq:two-layer-alpha} можно записать как однослойную~\eqref{ellipsoid-eq:alpha-homog} с внешним коэффициентом $L_j^{(2)}$ и \emph{эффективной} проницаемостью $\varepsilon_j^{\mathrm{eff}}$ ядра-в-оболочке, определяемой условием $\alpha_j^{(2)}=V_2(\varepsilon_j^{\mathrm{eff}}-\varepsilon_m)/[\varepsilon_m+L_j^{(2)}(\varepsilon_j^{\mathrm{eff}}-\varepsilon_m)]$; полученное $\varepsilon_j^{\mathrm{eff}}$ затем подставляют как «частицу» для следующей оболочки.

Для произвольного числа $N$ конфокальных слоёв процедура записывается как упорядоченное произведение $2\times2$ матриц переноса по каждой оси $j$ независимо. Введя послойный потенциальный «вектор состояния» $\bm{v}_p^{(j)}=(A_p,\,B_p)^{\mathsf T}$, имеем
\begin{equation}
\bm{v}_{p}^{(j)}=\bm{T}_p^{(j)}\,\bm{v}_{p-1}^{(j)},\qquad
\bm{v}_{N}^{(j)}=\Bigl(\textstyle\prod_{p=N}^{1}\bm{T}_p^{(j)}\Bigr)\,\bm{v}_{0}^{(j)},
\label{ellipsoid-eq:transfer-product}
\end{equation}
где матрица одной конфокальной границы $\xi_p$ между средами $\varepsilon_p$ (внутри) и $\varepsilon_{p+1}$ (снаружи), полученная из сшивки $\phi$ и $\varepsilon\,\partial_\xi\phi$ на гармониках $E_1(\xi),F_1(\xi)$, имеет структуру \cite{sihvola1990}
\begin{equation}
\bm{T}_p^{(j)}=\frac{1}{W_p}
\begin{pmatrix}
\varepsilon_{p+1}\,d_{F}-\varepsilon_{p}\,d'_{F} & \varepsilon_{p}\,d'_{F}-\varepsilon_{p+1}\,d_{F} \\[2pt]
\varepsilon_{p}\,d'_{E}-\varepsilon_{p+1}\,d_{E} & \varepsilon_{p+1}\,d_{E}-\varepsilon_{p}\,d'_{E}
\end{pmatrix},
\label{ellipsoid-eq:transfer-matrix}
\end{equation}
с $d_{E}=E_1(\xi_p)$, $d_{F}=F_1(\xi_p)$, $d'_{E}=\partial_\xi E_1|_{\xi_p}$, $d'_{F}=\partial_\xi F_1|_{\xi_p}$ и вронскианной нормировкой $W_p=\varepsilon_{p+1}(E_1 F'_1-E'_1 F_1)|_{\xi_p}$. Требование регулярности внутри ядра убирает убывающую гармонику ($B_0=0$), а связь излучённой ($B_N$) и приложенной ($A_N\!\propto\!E_0$) амплитуд на внешней границе даёт \emph{эффективную поляризуемость} вдоль оси $j$ в форме~\eqref{ellipsoid-eq:alpha-homog} с внешним коэффициентом $L_j^{(N)}$ и эффективной проницаемостью $\varepsilon_j^{\mathrm{eff}}$ всей вложенной структуры:
\begin{equation}
\alpha^{\mathrm{eff}}_{j}=V_N\,
\frac{\varepsilon^{\mathrm{eff}}_{j}-\varepsilon_m}
     {\varepsilon_m+L_j^{(N)}\bigl(\varepsilon^{\mathrm{eff}}_{j}-\varepsilon_m\bigr)} ,\qquad
V_N=\frac{4}{3}\pi a_N b_N c_N .
\label{ellipsoid-eq:alpha-eff}
\end{equation}
Эта факторизация --- ровно эллипсоидальный аналог скалярной импедансной рекурсии Адена--Керкера для сферы: вместо отношения Риккати--Бесселя на каждой сферической границе здесь по каждой главной оси перемножаются $2\times2$ матрицы~\eqref{ellipsoid-eq:transfer-matrix} на конфокальных эллипсоидальных границах \cite{sihvola1990,aden1951}.

\subsection{Сечения рассеяния и поглощения}

С эффективными главными поляризуемостями $\alpha^{\mathrm{eff}}_j$ (размерность объёма) частица --- точечный диполь $\bm{p}=\varepsilon_0\varepsilon_m\,\bm{\alpha}^{\mathrm{eff}}\!\cdot\!\bm{E}_0$, и стандартные рэлеевские сечения для падения, поляризованного вдоль оси $j$, суть \cite{bohren1983}
\begin{equation}
C_{\mathrm{sca},j}=\frac{\kappa^4}{6\pi}\,\bigl|\alpha^{\mathrm{eff}}_{j}\bigr|^2,
\qquad
C_{\mathrm{abs},j}=\kappa\,\mathrm{Im}\,\alpha^{\mathrm{eff}}_{j},
\qquad
C_{\mathrm{ext},j}=C_{\mathrm{abs},j}+C_{\mathrm{sca},j},
\label{ellipsoid-eq:rayleigh-cross}
\end{equation}
где $\kappa=\omega\sqrt{\varepsilon_m}/c$ --- волновое число во внешней среде. В дипольном пределе $C_{\mathrm{abs}}\propto\kappa$ доминирует над $C_{\mathrm{sca}}\propto\kappa^4$ для поглощающих частиц. Для случайно ориентированного ансамбля усредняем по трём осям,
\begin{equation}
\langle C_{\mathrm{ext}}\rangle=\tfrac13\bigl(C_{\mathrm{ext},a}+C_{\mathrm{ext},b}+C_{\mathrm{ext},c}\bigr),
\label{ellipsoid-eq:orientation-average}
\end{equation}
и аналогично для $C_{\mathrm{sca}},C_{\mathrm{abs}}$; в частности $\langle C_{\mathrm{sca}}\rangle=\dfrac{\kappa^4}{18\pi}\bigl(|\alpha^{\mathrm{eff}}_a|^2+|\alpha^{\mathrm{eff}}_b|^2+|\alpha^{\mathrm{eff}}_c|^2\bigr)$ \cite{bohren1983}. Три различных $L_j$ (и, в многослойном случае, три различных $\alpha^{\mathrm{eff}}_j$) проявляются как три раздельных резонансных пика.

\subsection{Область применимости и связь с библиотекой}

Сформулированный результат корректен строго при: (i) конфокальности всех границ~\eqref{ellipsoid-eq:confocal-req}; (ii) квазистатике $x=\kappa a\ll1$ (на практике $x\lesssim0.2$\,--\,$0.3$), когда запаздывание и высшие мультиполи $n\ge2$ пренебрежимы; (iii) линейных изотропных слоях. Вне квазистатики для триаксиального эллипсоида аналитического модуля нет (см.~\ref{ellipsoid-sec:ellipsoid}, п.~о невозможности векторного разделения): здесь библиотека ограничивается рэлеевской поляризуемостью~\eqref{ellipsoid-eq:alpha-eff}, которую при необходимости можно подать как дипольный вклад в Т-матричную суперпозицию (GMM) наряду с полноволновыми Т-матрицами сферы и сфероида. Таким образом, эллипсоид присутствует в семействе примитивов на честном уровне точности --- как точный диполь конфокальной слоистой структуры, а не как полноволновой решатель.

\section{Слоистый бесконечный цилиндр: метод тензорных функций Грина (ТФГ)}
\label{laycyl-sec:layered-cylinder}

Радиально-слоистый бесконечный круговой цилиндр (слои $i=1\ldots N$ с границами $a_1<\dots<a_N$, относительные $\varepsilon_i,\mu_i$; внешняя среда --- вакуум) допускает \emph{строгое} аналитическое решение методом тензорных функций Грина в форме эквивалентных линий передачи \cite{daylis2017,mitelman2013}. Задача разделяется в цилиндрических координатах, а многослойность сводится к послойной рекурсии матриц переноса --- цилиндрическому аналогу импедансно-адмитансной рекурсии слоистой сферы.

\subsection{Модальное разложение}
\label{laycyl-ssec:modal}
Поле представляется рядом Фурье по азимуту $n$ и (для наклонного падения) гармоникой вдоль оси:
\begin{equation}
E_z,\,H_z \;\sim\; \sum_{n=-\infty}^{\infty} [\,\cdot\,]\,e^{i(n\varphi + k_z z)},\qquad k_z=k_0\cos\theta_i,
\label{laycyl-eq:expansion}
\end{equation}
где $\theta_i$ --- угол падения, отсчитанный от оси цилиндра ($\theta_i=\pi/2$ --- нормальное падение). В слое $i$ продольные компоненты удовлетворяют уравнению Бесселя с радиальным волновым числом
\begin{equation}
\gamma_i=\sqrt{\varepsilon_i\mu_i\,k_0^2-k_z^2}=k_0\sqrt{\varepsilon_i\mu_i-\cos^2\theta_i}
\label{laycyl-eq:gamma}
\end{equation}
и раскладываются в устойчивом базисе $\{J_n(\gamma_i\rho),\,H_n^{(1)}(\gamma_i\rho)\}$ (функция Ханкеля первого рода --- исходящая при $e^{-i\omega t}$; базис корректен и при комплексных $\varepsilon_i$, т.\,е. для поглощающих слоёв). В сердцевине ($i=1$) присутствует только $J_n$ (регулярность на оси).

\subsection{Нормальное падение: расцепление TM/TE}
\label{laycyl-ssec:normal}
При $k_z=0$ поперечные компоненты $E_\varphi,H_\varphi$ не перемешивают $E_z$ и $H_z$: задача распадается на две независимые скалярные --- TM ($E\parallel$ ось) и TE ($H\parallel$ ось). Непрерывность продольной компоненты и тангенциального поля на каждой границе $\rho=a_j$ даёт $2\times2$ матрицу переноса; множитель в условии непрерывности
\begin{equation}
\eta_i^{\mathrm{TM}}=\frac{\sqrt{\varepsilon_i\mu_i}}{\mu_i},\qquad
\eta_i^{\mathrm{TE}}=\frac{\sqrt{\varepsilon_i\mu_i}}{\varepsilon_i}.
\label{laycyl-eq:eta}
\end{equation}
Прогон амплитуд от регулярной сердцевины наружу и сшивка с вакуумом дают поверхностную проводимость $G_n$ и коэффициент рассеяния
\begin{equation}
a_n=\frac{J_n'(x)-G_n\,J_n(x)}{H_n^{(1)\prime}(x)-G_n\,H_n^{(1)}(x)},\qquad x=k_0 a_N.
\label{laycyl-eq:an}
\end{equation}
При $N=1$ выражение~\eqref{laycyl-eq:an} тождественно сводится к коэффициентам однородного цилиндра Бора--Хаффмана \cite{bohren1983} (случай~I --- TM, случай~II --- TE).

\subsection{Наклонное падение: связь поляризаций}
\label{laycyl-ssec:oblique}
При $k_z\neq0$ граничные условия связывают E- и H-поляризации через член $\propto n\,k_z/\rho$ в $E_\varphi,H_\varphi$ (в терминологии эквивалентных линий Шабунина --- восьмиполюсник границы с коэффициентом связи $N_i\propto m\,h$, обращающимся в нуль при $h=0$ или $n=0$ \cite{daylis2017}). Поле в слое описывается уже четырьмя амплитудами ($E_z$: $J_n,H_n^{(1)}$; $H_z$: $J_n,H_n^{(1)}$), а непрерывность четырёх тангенциальных компонент $(E_z,H_z,E_\varphi,H_\varphi)$ даёт $4\times4$ матрицу переноса на слой. Для TM-падения (вертикальная поляризация) решение связанной системы на внешней границе даёт со-поляризационный коэффициент $A_n^{E}$ (рассеянный $E_z$) и кросс-поляризационный $A_n^{H}$ (рассеянный $H_z$); последний обращается в нуль при $\theta_i\to\pi/2$.

\subsection{Сечения и верификация}
\label{laycyl-ssec:cross}
Эффективности (нормировка на диаметр $2a$ и на падающую компоненту $E_z$):
\begin{equation}
Q_{\mathrm{sca}}=\frac{2}{x}\Big[\,|A_0^{E}|^2+|A_0^{H}|^2+2\!\sum_{n\ge1}\!\big(|A_n^{E}|^2+|A_n^{H}|^2\big)\Big],\quad
Q_{\mathrm{ext}}=-\frac{2}{x}\,\mathrm{Re}\Big[A_0^{E}+2\!\sum_{n\ge1}\!A_n^{E}\Big],
\label{laycyl-eq:cross}
\end{equation}
экстинкция --- по со-поляризации (оптическая теорема), $Q_{\mathrm{abs}}=Q_{\mathrm{ext}}-Q_{\mathrm{sca}}$. Реализация (\texttt{green\_tensor/cylinder.py}: \texttt{layered\_coeff}, \texttt{layered\_coeff\_oblique}, \texttt{cross\_sections\_layered/\_oblique}) проверена до машинной точности: (i)~$N=1$, нормальное падение --- против Бора--Хаффмана; (ii)~$N=1$, наклонное --- против коэффициентов одиночного цилиндра Уэйта \cite{wait1955} в форме Кавакліоğлу \cite{kavaklioglu2008} ($|\Delta|\sim10^{-16}$); (iii)~предельный переход $\theta_i\to\pi/2$; (iv)~инвариантность к подразбиению однородного цилиндра на одинаковые слои; (v)~сохранение энергии $Q_{\mathrm{abs}}\approx0$ для непоглощающих структур на всех углах.

\section{Конечный цилиндр: полуаналитическое решение методом сшивания мод}
\label{cylinder-cyl-sec:cylinder}

Конечный круговой цилиндр (радиус $a$, высота $L$, в общем случае многослойный по радиусу) --- первый из решателей \emph{уровня~B} (полуаналитика). В отличие от сферы и конфокального сфероида, для которых полноволновая векторная задача разделяется в \emph{одной} системе координат и приводит к замкнутой (диагональной или блочной по $m$) $T$-матрице, конечный цилиндр такого решения не допускает. Тем не менее он не требует пространственной сетки: поле в каждой подобласти раскладывается по \emph{замкнутым} собственным функциям, а конечность тела учитывается \emph{сшиванием мод} (mode matching) на стыке двух координатных поверхностей, дополненным условием на ребре. Результат --- усечённая связанная линейная система по модам, эквивалентная построению $T$-матрицы методом нуль-поля (EBCM) Уотермана \cite{waterman1971}; полученная $T$-матрица недиагональна в сферическом базисе ВСВФ и встраивается в схему GMM наравне с прочими примитивами.

\subsection{Почему конечный цилиндр не разделяется одним SOV}
\label{cylinder-cyl-ssec:nonsep}

Бесконечный круговой цилиндр разделяется в круговых цилиндрических координатах $(\rho,\varphi,z)$: скалярное уравнение Гельмгольца $(\nabla^2+k^2)\psi=0$ распадается на бесселево радиальное уравнение, гармоническую азимутальную зависимость $e^{im\varphi}$ ($m\in\mathbb{Z}$) и плоскую осевую $e^{ik_z z}$ \cite{stratton1941,eisenhart1934,bowman1987}. Конечный цилиндр ограничен \emph{двумя} разными координатными поверхностями --- боковой стенкой $\rho=a$ ($\rho=\mathrm{const}$) и двумя торцами $z=\pm L/2$ ($z=\mathrm{const}$), --- которые пересекаются по двум рёбрам $\rho=a,\ z=\pm L/2$. По теореме Эйзенхарта о штеккелевых системах \cite{eisenhart1934} полная сепарабельность скалярного уравнения требует, чтобы тело было ограничено единственной координатной поверхностью одной разделяющей системы (как сфера или сфероид); тело, ограниченное двумя различными координатными поверхностями с ребром между ними, не сводится к одному разделению переменных. Поэтому единого замкнутого ряда (как ряд Ми) для конечного цилиндра не существует: связь азимутального и осевого спектров на торцах и на боковой стенке порождает \emph{связанную} систему мод, а особенность поля у острого ребра требует отдельного условия (Мейкснера). Это и есть переход от строгой аналитики (уровень~A) к полуаналитике (уровень~B): базис остаётся аналитическим, но коэффициенты определяются из усечённой линейной системы, а не поканально.

\subsection{Собственные функции кругового цилиндра}
\label{cylinder-cyl-ssec:eigen}

Принята временн\'ая конвенция $e^{-i\omega t}$. В однородном слое с волновым числом $k=\omega\sqrt{\varepsilon\mu}$ (в общем случае комплексным) скалярные порождающие функции суть
\begin{equation}
\psi_{m}^{(J)}(\rho,\varphi,z;k_z)=\mathcal{Z}_m^{(J)}(k_\rho\rho)\,e^{im\varphi}\,e^{ik_z z},
\qquad
k_\rho^2+k_z^2=k^2,
\label{cylinder-cyl-eq:scalar-gen}
\end{equation}
где $m\in\mathbb{Z}$ --- азимутальный порядок, $k_z$ --- осевое волновое число, $k_\rho=\sqrt{k^2-k_z^2}$ --- поперечное, а $\mathcal{Z}_m^{(J)}$ --- цилиндрическая функция: $\mathcal{Z}_m^{(1)}=J_m$ (бесселева, регулярна на оси $\rho=0$) и $\mathcal{Z}_m^{(3)}=H_m^{(1)}$ (ханкелева первого рода) \cite{stratton1941,bowman1987,dlmf}. Выбор $H_m^{(1)}(k_\rho\rho)\sim\sqrt{2/(\pi k_\rho\rho)}\,e^{i(k_\rho\rho-m\pi/2-\pi/4)}$ при $\rho\to\infty$ задаёт волну, \emph{уходящую} от оси, в согласии с конвенцией $e^{-i\omega t}$ --- прямой цилиндрический аналог сферического $h_n^{(1)}$ \cite{stratton1941,dlmf}. Ветвь $k_\rho$ фиксируется условием $\operatorname{Im}k_\rho\ge0$ (затухание/уход в дальней зоне); при $k_z^2>k^2$ ($k_\rho$ чисто мнимое) моды эванесцентны по $\rho$ и описываются модифицированными функциями $I_m,K_m$.

\paragraph{Векторные цилиндрические волновые функции.} По образцу Стрэттона \cite{stratton1941} соленоидальные решения векторного уравнения $\nabla\times\nabla\times\bm{E}-k^2\bm{E}=\bm0$ строятся через пилотный вектор $\hat{\bm z}$ (ось цилиндра):
\begin{equation}
\bm{m}_{m}^{(J)}(k_z)=\nabla\times\!\big(\hat{\bm z}\,\psi_{m}^{(J)}\big),
\qquad
\bm{n}_{m}^{(J)}(k_z)=\frac{1}{k}\,\nabla\times\bm{m}_{m}^{(J)}(k_z),
\label{cylinder-cyl-eq:mn-def}
\end{equation}
причём $\bm{m}_{m}^{(J)}$ не имеет $z$-компоненты (поперечно-электрический, TE, относительно оси), а $\bm{n}_{m}^{(J)}$ несёт продольную компоненту $E_z\propto(k_\rho^2/k)\,\psi_m^{(J)}$ (поперечно-магнитный, TM). Третье (продольное) семейство $\bm{l}_m^{(J)}=\nabla\psi_m^{(J)}$ безвихревое и в бездивергентном поле не используется. В явном виде
\begin{align}
\bm{m}_{m}^{(J)}&=\Big[\tfrac{im}{\rho}\,\mathcal{Z}_m^{(J)}(k_\rho\rho)\,\hat{\bm\rho}
-k_\rho\,\mathcal{Z}_m^{(J)\prime}(k_\rho\rho)\,\hat{\bm\varphi}\Big]e^{im\varphi+ik_z z},
\label{cylinder-cyl-eq:m-explicit}\\
\bm{n}_{m}^{(J)}&=\frac{1}{k}\Big[ik_z k_\rho\,\mathcal{Z}_m^{(J)\prime}\,\hat{\bm\rho}
-\tfrac{m k_z}{\rho}\,\mathcal{Z}_m^{(J)}\,\hat{\bm\varphi}
+k_\rho^2\,\mathcal{Z}_m^{(J)}\,\hat{\bm z}\Big]e^{im\varphi+ik_z z}.
\label{cylinder-cyl-eq:n-explicit}
\end{align}
Азимутальная зависимость $e^{im\varphi}$ \emph{развязывает} разные порядки $m$ (осевая симметрия тела), но при наклонном падении (произвольный $k_z\ne0$) граничные условия на боковой стенке смешивают TE и TM: чистое расщепление $\bm m\!\leftrightarrow\!\bm n$ имеет место лишь при нормальном падении ($k_z=0$, поле не зависит от $z$), как и в классической 2D-задаче о бесконечном цилиндре \cite{bowman1987,stratton1941}. Это фундаментальная черта цилиндра, отличающая его от сферы.

\paragraph{Многослойность (радиальное расслоение).} Коаксиальные слои $\rho=a_1<\dots<a_S=a$ обрабатываются \emph{до} торцевого сшивания: при фиксированных $(m,k_z)$ непрерывность тангенциальных $E_\varphi,E_z,H_\varphi,H_z$ на каждой цилиндрической границе $\rho=a_s$ связывает амплитуды соседних слоёв. Поскольку при $k_z\ne0$ TE и TM сцеплены, передача идёт $4\times4$ матрицей по паре $(\bm m,\bm n)$ (для $k_z=0$ она распадается на два скалярных $2\times2$ канала). Это прямой цилиндрический аналог рекурсии Адена--Керкера: накопленная по слоям радиальная матрица переноса $\bm{R}^{(m)}(k_z)$ при заданном $(m,k_z)$ связывает поле на оси (регулярность $J_m$ в ядре, отсутствие $H_m^{(1)}$) с полем на боковой поверхности $\rho=a$, и её размер не зависит от числа слоёв \cite{stratton1941,bohren1983}. После этого шага цилиндр любой радиальной структуры заменяется эффективным импедансным граничным оператором на $\rho=a$, и дальнейшее сшивание мод не зависит от деталей слоёв.

\subsection{Разложение по подобластям и сшивание мод}
\label{cylinder-cyl-ssec:matching}

Конечность тела вводится сопоставлением \emph{внутреннего} (или приповерхностного) модального поля \emph{внешнему} мультипольному полю через всю замкнутую границу --- боковую стенку и оба торца. Внешнее поле естественно записать сразу в сферических ВСВФ с общим началом в центре цилиндра (это и есть базис, нужный для GMM):
\begin{align}
\bm{E}^{\mathrm{inc}}(\bm r)&=\sum_{n=1}^{\infty}\sum_{\mu=-n}^{n}
\Big[d_{\mu n}^{(1)}\,\mathrm{Rg}\,\bm{M}_{\mu n}(\bm r)+d_{\mu n}^{(2)}\,\mathrm{Rg}\,\bm{N}_{\mu n}(\bm r)\Big],
\label{cylinder-cyl-eq:Einc}\\
\bm{E}^{\mathrm{sca}}(\bm r)&=\sum_{n=1}^{\infty}\sum_{\mu=-n}^{n}
\Big[c_{\mu n}^{(1)}\,\bm{M}_{\mu n}(\bm r)+c_{\mu n}^{(2)}\,\bm{N}_{\mu n}(\bm r)\Big],
\label{cylinder-cyl-eq:Esca}
\end{align}
с регулярными ($J=1$, $j_n$) и исходящими ($J=3$, $h_n^{(1)}$) ВСВФ \cite{stratton1941,mishchenko2002}. Внутреннее поле (после радиальной редукции слоёв) записывается как непрерывный по $k_z$ спектр цилиндрических мод \eqref{cylinder-cyl-eq:m-explicit}--\eqref{cylinder-cyl-eq:n-explicit}, регулярных на оси:
\begin{equation}
\bm{E}^{\mathrm{int}}(\bm r)=\sum_{m}\int_{-\infty}^{\infty}\!\!\Big[\alpha_m(k_z)\,\bm{m}_m^{(1)}(k_z)+\beta_m(k_z)\,\bm{n}_m^{(1)}(k_z)\Big]dk_z .
\label{cylinder-cyl-eq:Eint}
\end{equation}
(Осевая ограниченность $|z|\le L/2$ означает, что спектр по $k_z$ непрерывен и на торцах раскладывается по поперечным модам кругового сечения $J_m(\gamma_{mp}\rho/a)$.) На обеих координатных поверхностях требуется непрерывность тангенциальных $\bm E$ и $\bm H$:
\begin{equation}
\hat{\bm n}\times\bm{E}^{\mathrm{out}}=\hat{\bm n}\times\bm{E}^{\mathrm{int}},\qquad
\hat{\bm n}\times\bm{H}^{\mathrm{out}}=\hat{\bm n}\times\bm{H}^{\mathrm{int}}
\quad\text{на }\ \rho=a\ (\hat{\bm n}=\hat{\bm\rho})\ \text{и}\ z=\pm L/2\ (\hat{\bm n}=\pm\hat{\bm z}),
\label{cylinder-cyl-eq:bc}
\end{equation}
где $\bm{E}^{\mathrm{out}}=\bm{E}^{\mathrm{inc}}+\bm{E}^{\mathrm{sca}}$. Поскольку поверхности разные, проекции базисов не ортогональны: чтобы спроецировать условия \eqref{cylinder-cyl-eq:bc} на единый счётный набор, сферические ВСВФ переразлагают на торцах по дисковым модам Бесселя $\{J_m(\gamma_{mp}\rho/a)e^{im\varphi}\}_{p}$, а на боковой стенке --- по фурье-модам $\{e^{im\varphi}e^{i\pi q z/L}\}_q$ (либо непосредственно по $\int dk_z$). Это даёт связь модальных коэффициентов через \emph{интегралы перекрытия} между цилиндрическими и сферическими гармониками (по сути присоединённые функции Лежандра, интегрированные с бесселевыми ядрами по торцам и стенке) --- цилиндрический аналог матриц перекрытия $G^{(m)}_{ll'}$ сфероидальной задачи. После усечения ($|m|\le m_{\max}$, $n\le n_{\max}$, $p,q\le P$) получается \emph{связанная} линейная система
\begin{equation}
\bm{Q}^{(m)}\,
\begin{pmatrix}\bm{\alpha}_m\\ \bm{\beta}_m\\ \bm{c}_m\end{pmatrix}
=\bm{P}^{(m)}\,\bm{d}_m ,
\label{cylinder-cyl-eq:modal-system}
\end{equation}
где блоки $\bm{Q}^{(m)},\bm{P}^{(m)}$ собраны из значений цилиндрических и сферических функций на границах, их нормальных производных и интегралов перекрытия. Развязка по $m$ сохраняется (тело осесимметрично), а порядки $n$ и каналы TE/TM связываются именно через торцевое/боковое сшивание --- это и есть источник недиагональности $T$. Исключение внутренних амплитуд $(\bm\alpha_m,\bm\beta_m)$ даёт связь $\bm{c}_m=\bm{T}^{(m)}\bm{d}_m$.

\subsection{Условие на ребре (Мейкснера) и сходимость}
\label{cylinder-cyl-ssec:edge}

Острые рёбра $\rho=a,\ z=\pm L/2$ --- источник особенности поля. Корректная (единственная) постановка требует наложения \emph{условия Мейкснера на ребре}: плотность энергии интегрируема в любой окрестности ребра, а компоненты поля ведут себя как $\sim\delta^{\,\tau-1}$ при расстоянии $\delta\to0$ до ребра, где показатель $\tau$ определяется углом раскрыва. Для прямого ($90^\circ$) ребра идеального проводника (внешний угол раскрыва диэлектрика $\Phi=3\pi/2$, $\tau=\pi/\Phi=2/3$) тангенциальное к ребру поле конечно, а нормальные компоненты $E_\rho,E_z$ и плотность тока расходятся как $\delta^{\tau-1}=\delta^{-1/3}$ (диэлектрическое ребро даёт другой, зависящий от контраста, нецелый показатель) \cite{bowman1987,mishchenko2002}. Физически условие на ребре исключает нефизические решения (паразитные токи на ребре) и тем самым замыкает систему \eqref{cylinder-cyl-eq:modal-system}; численно оно реализуется выбором базиса с правильной асимптотикой у ребра либо взвешиванием невязки, что резко ускоряет сходимость усечённого ряда. Без учёта особенности модальный ряд сходится лишь алгебраически и обнаруживает явление Гиббса у рёбер. Поэтому скорость сходимости по $(n_{\max},P)$ управляется не только параметром размера $ka$, но и \emph{аспектным отношением} $L/(2a)$ и остротой ребра: тонкие диски ($L\ll a$) и длинные стержни ($L\gg a$) требуют существенно б\'ольшего усечения и тщательного учёта реберной асимптотики. Это типичная для уровня~B честная оговорка: решение точно в пределе $n_{\max}\to\infty$, но матрица связана, а сходимость зависит от формы.

\subsection{Маршрут к сферической $T$-матрице и связь с EBCM}
\label{cylinder-cyl-ssec:tmatrix}

Система \eqref{cylinder-cyl-eq:modal-system} уже сформулирована в сферическом базисе на входе ($\bm d_m$) и выходе ($\bm c_m$), поэтому исключение внутренних амплитуд прямо даёт блочную по $m$ $T$-матрицу
\begin{equation}
\bm{c}_m=\bm{T}^{(m)}\,\bm{d}_m,\qquad
\bm{T}=\bigoplus_{m}\bm{T}^{(m)},\qquad
\bm{T}^{(m)}=\begin{pmatrix}\bm{T}^{11}_{(m)}&\bm{T}^{12}_{(m)}\\[2pt]\bm{T}^{21}_{(m)}&\bm{T}^{22}_{(m)}\end{pmatrix},
\label{cylinder-cyl-eq:Tblock}
\end{equation}
недиагональную по $n$ и смешивающую TE/TM ($\bm{T}^{12},\bm{T}^{21}\ne\bm0$ при наклонном падении) --- в полном соответствии с \eqref{tmatrix-eq:tm-Td} и в отличие от диагональной $\bm{T}=\operatorname{diag}(-a_n,-b_n)$ сферы. Эквивалентный и численно предпочтительный путь --- метод нуль-поля (EBCM) Уотермана \cite{waterman1971,mishchenko2002}: $T$-матрица строится как
\begin{equation}
\bm{T}=-\,\mathrm{Rg}\bm{Q}\;\bm{Q}^{-1},
\label{cylinder-cyl-eq:ebcm}
\end{equation}
где матрицы $\bm{Q}$ и $\mathrm{Rg}\bm{Q}$ суть поверхностные интегралы по \emph{всей} границе цилиндра $\Sigma=\{\rho=a\}\cup\{z=\pm L/2\}$ от произведений ВСВФ ($\bm{N}^{(3,1)},\bm{M}^{(3,1)}$ и регулярных $\mathrm{Rg}$-аналогов) с тангенциальными поверхностными полями, выражаемыми через тот же набор цилиндрических собственных функций. Острые рёбра в EBCM учитываются дискретизацией контура с правильной реберной асимптотикой; для осесимметричного тела интеграл по $\varphi$ берётся аналитически, оставляя одномерный контурный интеграл по образующей --- отсюда отсутствие пространственной сетки \cite{waterman1971,mishchenko2002}. Оба маршрута --- прямое сшивание мод и EBCM --- дают одну и ту же $\bm{T}^{(m)}$; их совпадение служит внутренней проверкой. Полученная матрица переразлагается/используется как есть в сферическом базисе и подаётся в аддитивные теоремы Стейна--Крузана и самосогласованную систему GMM \eqref{tmatrix-eq:tm-linsys} наравне с прочими примитивами; при произвольной ориентации применяется поворот Вигнера \eqref{tmatrix-eq:tm-rot}. Концептуально это та же стратегия «непрерывного базиса с перепроекцией на сферические ВСВФ», что и в обобщённом методе разделения переменных / EBCM с неклассическим базисом \cite{farafonov2012}.

\subsection{Дальнее поле и сечения}
\label{cylinder-cyl-ssec:farfield}

Из исходящих коэффициентов $\bm c=\bm T\bm d$ дальнее поле и сечения вычисляются по тем же формулам, что и для прочих примитивов (раздел~\ref{tmatrix-sec:tm-farfield}). В дальней зоне
\begin{equation}
\bm{E}^{\mathrm{sca}}\xrightarrow[kr\to\infty]{}\frac{e^{ikr}}{kr}\,\bm{F}(\hat{\bm r}),\qquad
\bm{F}(\hat{\bm r})=\sum_{n,\mu}\big[(-i)^{\,n+1}c_{\mu n}^{(1)}\bm{\mathcal X}^{M}_{\mu n}(\hat{\bm r})+(-i)^{\,n}c_{\mu n}^{(2)}\bm{\mathcal X}^{N}_{\mu n}(\hat{\bm r})\big],
\label{cylinder-cyl-eq:far}
\end{equation}
с поперечными угловыми гармониками $\bm{\mathcal X}_{\mu n}$ \cite{mishchenko2002,bohren1983}; раздельные фазовые множители $(-i)^{n+1}$ (для $\bm M$-семейства) и $(-i)^{n}$ (для $\bm N$-семейства) суть прямое следствие дальнезонной асимптотики $h_n^{(1)}(kr)\to(-i)^{n+1}e^{ikr}/(kr)$ и производной комбинации $\big(kr\,h_n^{(1)}\big)'/(kr)\to(-i)^{n}e^{ikr}/(kr)$. Бистатическое сечение (RCS) есть $d\sigma/d\Omega=|\bm F(\hat{\bm r})|^2/k^2$ (согласовано с нормировкой $\bm E^{\mathrm{sca}}\to\bm F\,e^{ikr}/(kr)$ выше). Полные сечения --- стандартно через след $T$-матрицы (как в \eqref{spheroid-eq:cross}) либо коэффициентно через оптическую теорему \eqref{tmatrix-eq:tm-Cext}:
\begin{equation}
C_{\mathrm{ext}}=-\frac{1}{k^2}\operatorname{Re}\sum_{n,\mu}\big[d_{\mu n}^{(1)}(c_{\mu n}^{(1)})^{*}+d_{\mu n}^{(2)}(c_{\mu n}^{(2)})^{*}\big],
\qquad
C_{\mathrm{sca}}=\frac{1}{k^2}\sum_{n,\mu}\big(|c_{\mu n}^{(1)}|^2+|c_{\mu n}^{(2)}|^2\big),
\label{cylinder-cyl-eq:cross}
\end{equation}
причём $C_{\mathrm{abs}}=C_{\mathrm{ext}}-C_{\mathrm{sca}}\ge0$ для поглощающих слоёв \cite{bohren1983,mishchenko2002}. Унитарность $\bm{T}^{\dagger}\bm{T}+\operatorname{Re}\bm{T}=\bm0$ для непоглощающего цилиндра \cite{waterman1971} служит сквозным индикатором сходимости усечения и корректности учёта ребра.

\subsection{Предельные переходы и контроль}
\label{cylinder-cyl-ssec:limits}

\paragraph{Бесконечный цилиндр ($L\to\infty$).} При неограниченном удлинении торцы уходят на бесконечность, торцевое сшивание исчезает, и осевой спектр коллапсирует к фиксированному $k_z$ падающей волны (трансляционная инвариантность вдоль $z$). Система \eqref{cylinder-cyl-eq:modal-system} распадается по $m$ на независимые задачи о бесконечном цилиндре, и для каждого $m$ восстанавливается \emph{точный} классический 2D-ряд: коэффициенты отражения выражаются через кросс-произведения $J_m(k_\rho^{\mathrm{int}}a),H_m^{(1)}(k_\rho^{\mathrm{ext}}a)$ и их производные, причём при нормальном падении ($k_z=0$) TE и TM полностью развязываются \cite{bowman1987,stratton1941}. Это даёт прямой аналитический контроль: реализация конечного цилиндра при больших $L/a$ должна сходиться к табличному решению бесконечного цилиндра.

\paragraph{Прочие проверки.} (i)~Тонкий диск $L\ll a$ при $ka\ll1$ --- сверка с квазистатической поляризуемостью сплюснутого сфероида (предел уровня~A). (ii)~При $ka\ll1,\ kL\ll1$ --- рэлеевский предел и баланс энергии $C_{\mathrm{abs}}\ge0$. (iii)~Совпадение $\bm{T}^{(m)}$, полученной прямым сшиванием мод и методом нуль-поля \eqref{cylinder-cyl-eq:ebcm}. (iv)~Кросс-валидация сечений и диаграмм с независимым MoM/FEM. Честная оговорка уровня~B: $\bm{T}$ недиагональна, а сходимость по $(n_{\max},P)$ управляется $ka$, аспектным отношением $L/(2a)$ и реберной асимптотикой; для предельных форм (очень тонкие диски или длинные стержни) требуется повышенное усечение и аккуратный учёт условия Мейкснера \cite{bowman1987,mishchenko2002}.

\section{Конечный (усечённый) конус: полуаналитическое решение}
\label{cone-cone-sec:cone}

Конус (полный круговой, либо усечённый --- ограниченный снизу сферической шапкой или плоским основанием) принадлежит уже не к строгому семейству примитивов, ограниченных \emph{единственной} координатной поверхностью, а к классу \textbf{TIER-B} полуаналитических тел. Боковая поверхность кругового конуса есть координатная поверхность $\theta=\theta_0$ сферической (точнее --- сферо-конической) системы, и в ней скалярное уравнение Гельмгольца разделяется \cite{eisenhart1934,stratton1941}. Однако \emph{конечность} тела вводит вторую границу --- сферическую шапку $r=a$ (либо плоское основание), --- которая принадлежит другому координатному семейству и пересекает боковую поверхность по ребру. Тело, ограниченное \emph{двумя разными} координатными поверхностями, не решается одним разделением переменных; оно требует \emph{сшивки мод} (mode-matching) между сепарабельными подобластями с замкнутыми собственными функциями (без пространственной сетки), с явным учётом условия на ребре (Мейкснера). Эквивалентно $T$-матрица строится методом нуль-поля (EBCM Уотермена) на тех же собственных функциях \cite{waterman1971}. Честно отметим заранее: \emph{это самый трудный примитив библиотеки} --- машинерия присоединённых функций Лежандра \emph{нецелой} степени, слабая поддержка в стандартных библиотеках и связь порядков при сшивке делают его принципиально более тяжёлым, чем сфера или даже сфероид.

\subsection{Сферо-конические координаты и разделение для кругового конуса}
\label{cone-cone-ssec:sepvar}

Боковая поверхность кругового конуса с полууглом раствора $\theta_0$ совпадает с поверхностью $\theta=\theta_0$ обычной сферической системы $(r,\theta,\varphi)$ (для эллиптического конуса --- с координатной поверхностью сферо-конической системы, в которой угловая часть управляется уравнением Ламе) \cite{klinkenbusch2007,bowman1987}. Рассмотрим круговой случай. Скалярное уравнение Гельмгольца $(\nabla^2+k^2)\psi=0$ разделяется подстановкой
\begin{equation}
\psi_{\nu m}(r,\theta,\varphi)=z_\nu(kr)\,P_\nu^{m}(\cos\theta)\,e^{im\varphi},
\label{cone-cone-eq:scalar-gen}
\end{equation}
где $z_\nu$ --- сферическая функция Бесселя \emph{нецелого} индекса (радиальная зависимость $z_\nu(kr)=\sqrt{\pi/2kr}\,Z_{\nu+1/2}(kr)$, $Z=J,Y,H^{(1)}$), $m\in\mathbb{Z}$ --- азимутальный порядок (периодичность по $\varphi$), а $P_\nu^{m}(\cos\theta)$ --- присоединённая функция Лежандра \emph{нецелой степени} $\nu$ \cite{stratton1941,dlmf}. Радиальное и угловое уравнения суть
\begin{align}
\frac{1}{r^2}\frac{\dd}{\dd r}\!\left(r^2\frac{\dd z_\nu}{\dd r}\right)
&+\left(k^2-\frac{\nu(\nu+1)}{r^2}\right)z_\nu=0,
\label{cone-cone-eq:radial}\\[3pt]
\frac{1}{\sin\theta}\frac{\dd}{\dd\theta}\!\left(\sin\theta\,\frac{\dd P_\nu^m}{\dd\theta}\right)
&+\left(\nu(\nu+1)-\frac{m^2}{\sin^2\theta}\right)P_\nu^m=0,
\label{cone-cone-eq:angular}
\end{align}
с общей постоянной разделения $\nu(\nu+1)$. Ключевое отличие от сферы: на сфере регулярность \emph{на обоих} полюсах $\theta=0,\pi$ квантует степень в целые $\nu=l\ge|m|$; на конусе же область занимает только $0\le\theta\le\theta_0$ (внешность конуса) или $\theta_0\le\theta\le\pi$ (внутренность), \emph{один} полюс исключён конусом, а граничное условие на $\theta=\theta_0$ заменяет утраченное условие регулярности. Это снимает целочисленность и порождает \emph{дискретный спектр нецелых} степеней $\nu$ \cite{klinkenbusch2007,bowman1987}.

\paragraph{Характеристическое уравнение для степеней $\nu$.}
Для идеально проводящего (PEC) конуса тангенциальные компоненты строятся из двух дебаевских потенциалов; в скалярной формулировке потенциалы типа TM (электрический, ассоциированный с $\bm N$) обращают в нуль на боковой поверхности саму функцию, а типа TE (магнитный, ассоциированный с $\bm M$) --- её нормальную производную. Для угловой функции, регулярной на доступном полюсе, это даёт два набора собственных степеней $\{\nu_p^{\mathrm{TM}}\}$, $\{\nu_q^{\mathrm{TE}}\}$ как корни \cite{klinkenbusch2007,bowman1987,stratton1941}
\begin{align}
P_{\nu}^{m}(\cos\theta_0)&=0 \quad(\text{TM, условие Дирихле}),
\label{cone-cone-eq:eig-TM}\\
\left.\frac{\dd}{\dd\theta}P_{\nu}^{m}(\cos\theta)\right|_{\theta=\theta_0}&=0 \quad(\text{TE, условие Неймана}),
\label{cone-cone-eq:eig-TE}
\end{align}
по одному трансцендентному уравнению на каждое $m$. Для диэлектрического конуса вместо \eqref{cone-cone-eq:eig-TM}--\eqref{cone-cone-eq:eig-TE} ставится сшивка тангенциальных $\bm E,\bm H$ на $\theta=\theta_0$ между внутренними и внешними угловыми модами, что связывает соседние степени; в простейшей PEC-постановке степени $\nu$ определены \emph{только геометрией} $(\theta_0,m)$ и не зависят от размера. Угловые функции $\{P_{\nu_p}^m(\cos\theta)\}$ при фиксированном $m$ образуют полный ортогональный (с весом $\sin\theta$) базис на интервале $[\, \cdot\,,\theta_0]$ --- конический аналог $\{P_l^m\}$ сферы \cite{klinkenbusch2007,dlmf}.

\subsection{Точное решение полу-бесконечного конуса как строительный блок}
\label{cone-cone-ssec:greendyad}

Полу-бесконечный круговой (и эллиптический) конус --- это та подобласть, в которой задача решается \emph{строго аналитически}, и она служит несущим блоком для конечного тела. Клинкенбуш \cite{klinkenbusch2007} строит диадную (тензорную) функцию Грина конуса методом сферо-мультипольного разложения: она представляется билинейными произведениями сферических функций Бесселя радиального аргумента и угловых собственных функций нецелой степени (для эллиптического конуса --- периодических и непериодических функций Ламе). Схематически, с мультииндексом $\nu=(\nu,m)$ собственных мод \eqref{cone-cone-eq:eig-TM}--\eqref{cone-cone-eq:eig-TE},
\begin{equation}
\bar{\bm{G}}^{\mathrm{cone}}(\bm{r},\bm{r}')=
\sum_{\nu}\Big[\bm{M}_{\nu}^{(1)}(\bm{r}_<)\,\bm{M}_{\nu}^{(3)}(\bm{r}_>)
+\bm{N}_{\nu}^{(1)}(\bm{r}_<)\,\bm{N}_{\nu}^{(3)}(\bm{r}_>)\Big]/\mathcal{N}_\nu,
\label{cone-cone-eq:greendyad}
\end{equation}
где $\bm{M}_\nu,\bm{N}_\nu$ --- векторные конические волновые функции, построенные по \eqref{cone-cone-eq:scalar-gen} стандартным образом ($\bm M_\nu=\nabla\times(\bm r\,\psi_{\nu m}),\ \bm N_\nu=k^{-1}\nabla\times\bm M_\nu$), $\mathcal{N}_\nu$ --- нормы, $\bm r_{<,>}$ --- меньший/больший по $r$ аргумент. Поверхностный ток на PEC-конусе и рассеянное дальнее поле получаются сворачиванием \eqref{cone-cone-eq:greendyad} с падающим полем \cite{klinkenbusch2007}. Существенная численная тонкость: ряды по $\nu$ \emph{не} сходятся прямым суммированием (вершина конуса --- особенность поля), и для физически корректной асимптотики применяется линейное преобразование последовательностей (sequence transformation) \cite{klinkenbusch2007}. Этот блок задаёт «конический» отклик без учёта конечности; конечность добавляется сшивкой ниже.

\subsection{Конечность через сшивку мод (mode-matching)}
\label{cone-cone-ssec:modematch}

Усечём конус сферической шапкой $r=a$ (либо плоским основанием). Пространство делится на две сепарабельные подобласти: \emph{коническую} ($r<a$, $\theta\in[\,\cdot\,,\theta_0]$) с собственными модами нецелой степени $\nu$ из \eqref{cone-cone-eq:scalar-gen}--\eqref{cone-cone-eq:eig-TE} и \emph{внешнюю} ($r>a$, полный телесный угол) с обычными целочисленными сферическими мультиполями $\{\bm M_{nm}^{(J)},\bm N_{nm}^{(J)}\}$, $n=|m|,|m|{+}1,\dots$ \cite{stratton1941,mishchenko2002}. На разделяющей поверхности $r=a$ непрерывность тангенциальных $\bm E,\bm H$ сшивает два базиса. Поскольку при фиксированном $m$ азимутальная зависимость $e^{im\varphi}$ общая (порядки $m$ \emph{развязаны}), а угловые профили различны ($P_\nu^m$ против $P_n^m$ на разных $\theta$-интервалах), проекция одной системы на другую порождает прямоугольные \emph{матрицы перекрытия}
\begin{equation}
G^{(m)}_{n\nu}=\frac{\displaystyle\int_{\cos\theta_0}^{1} P_n^{m}(t)\,P_\nu^{m}(t)\,\dd t}
{\Big[\int P_\nu^{m\,2}\,\dd t\Big]^{1/2}\Big[\int P_n^{m\,2}\,\dd t\Big]^{1/2}},
\qquad t=\cos\theta,
\label{cone-cone-eq:overlap}
\end{equation}
которые связывают \emph{нецелые} $\nu$-моды конуса с \emph{целыми} $n$-мультиполями внешности. Сшивка тангенциальных компонент на $r=a$ (с радиальными функциями $z_\nu(ka)$ внутри и $j_n(ka),h_n^{(1)}(ka)$ снаружи, и их производными) приводит к \emph{связанной} (усечённой по $\nu\le\nu_{\max}$, $n\le n_{\max}$) линейной системе вида
\begin{equation}
\sum_{\nu}\!\Big[\bm{Q}^{(m)}_{\mathrm{int}}\Big]_{n\nu}\!c_\nu
=\sum_{n'}\!\Big[\bm{Q}^{(m)}_{\mathrm{ext}}\Big]_{nn'}\!\big(d_{n'}+f_{n'}\big),
\qquad
\bm f^{(m)}=\bm{T}^{(m)}\,\bm d^{(m)},
\label{cone-cone-eq:modematch-system}
\end{equation}
где $\bm{Q}^{(m)}_{\mathrm{int}},\bm{Q}^{(m)}_{\mathrm{ext}}$ --- блочные граничные матрицы (составленные из $G^{(m)}_{n\nu}$ \eqref{cone-cone-eq:overlap}, радиальных функций и их производных на $r=a$), $c_\nu$ --- амплитуды конических мод, $d_{n'},f_{n'}$ --- коэффициенты падающего и рассеянного сферических полей. Исключение внутренних амплитуд $c_\nu$ даёт искомую блочную (по $m$) $T$-матрицу. Именно \eqref{cone-cone-eq:modematch-system} --- источник недиагональности: \emph{целые} порядки $n$ оказываются связанными между собой через общий набор \emph{нецелых} $\nu$, и сшивка перемешивает TE/TM при наклонном падении (как и у сфероида, расщепление чисто TE/TM возможно лишь при осевом $m=0$).

\paragraph{Эквивалентность EBCM/нуль-поля.}
Альтернативно ту же $T$-матрицу можно построить методом нуль-поля Уотермена \cite{waterman1971}: поверхностные интегралы $\bm Q$- и $\mathrm{Rg}\,\bm Q$-матриц берутся по \emph{полной} границе усечённого конуса (боковая поверхность $\theta=\theta_0$ плюс шапка/основание), а $\bm{T}=-\mathrm{Rg}\,\bm{Q}\,\bm{Q}^{-1}$. На гладких участках границы поверхностные интегралы вычисляются в \emph{замкнутой} форме на введённых собственных функциях (без сетки), а ребро на стыке участков обрабатывается отдельно (см.\ ниже). Оба пути --- прямая сшивка \eqref{cone-cone-eq:modematch-system} и EBCM --- математически эквивалентны и дают одну и ту же недиагональную $\bm{T}^{(m)}$ \cite{waterman1971,mishchenko2002}.

\subsection{Условия на вершине и ребре основания (Мейкснер)}
\label{cone-cone-ssec:edge}

У усечённого конуса две сингулярные линии. \emph{Вершина} ($r\to0$) --- точка схождения боковой поверхности; вблизи неё \emph{потенциал} ведёт себя как $r^{\nu_{\min}}$, а сами компоненты поля (производные потенциала) --- как $r^{\nu_{\min}-1}$, с наименьшей собственной степенью $\nu_{\min}(\theta_0,m)$ из \eqref{cone-cone-eq:eig-TM}--\eqref{cone-cone-eq:eig-TE} \cite{klinkenbusch2007,bowman1987}. Сингулярность поля в вершине ($|\bm E|\sim r^{\nu_{\min}-1}\to\infty$) возникает лишь когда $\nu_{\min}<1$ --- то есть для конфигураций с острым \emph{проникающим} в среду проводящим шипом (большой раствор проводника); именно тогда ряды \eqref{cone-cone-eq:greendyad} сходятся медленно. Для тонкой иглы ($\theta_0\to0$) наоборот $\nu_{\min}$ велико и поле в вершине регулярно (ср.\ предельный переход $\theta_0\to0$ ниже). \emph{Ребро основания} (рим, где $\theta=\theta_0$ встречает $r=a$) --- линия пересечения двух границ; здесь поле подчиняется условию Мейкснера: компоненты, нормальные к ребру, могут расходиться, но не быстрее $\rho^{-1/2}$ ($\rho$ --- расстояние до ребра), а плотность энергии интегрируема, $\int |\bm E|^2\,\dd V<\infty$ \cite{bowman1987,klinkenbusch2007}. В дискретизованной системе \eqref{cone-cone-eq:modematch-system} условие Мейкснера реализуется выбором/взвешиванием базисных мод так, чтобы воспроизвести правильную сингулярность $\rho^{\tau}$ ($\tau$ зависит от двугранного угла ребра), и определяет скорость сходимости по $\nu_{\max},n_{\max}$. Игнорирование ребра ведёт к псевдо-осцилляциям (явление Гиббса) в поверхностном токе и к ложному замедлению сходимости; корректный учёт --- обязательное условие точности этого примитива \cite{bowman1987,mishchenko2002}.

\subsection{Многослойность: радиальная рекурсия внутри конуса}
\label{cone-cone-ssec:layers}

Стратификация поддерживается \emph{до} сшивки на ребре. Концентрические по $r$ слои (сферические оболочки внутри конического телесного угла, $r=a_1<a_2<\dots<a_{N_L}=a$, общая боковая граница $\theta=\theta_0$) сохраняют угловой базис $\{P_\nu^m(\cos\theta)\}$ неизменным (степени $\nu$ заданы геометрией $\theta_0$), а радиальную зависимость переносят послойно. Для каждой собственной моды $\nu$ радиальное уравнение \eqref{cone-cone-eq:radial} --- то же, что и для сферы, лишь с нецелым индексом, поэтому работает \emph{та же} дробно-линейная рекурсия импеданса/адмитанса Адена--Керкера, что и в каноническом ядре (см.\ \S\ref{sphere-sec:sphere}), с заменой целого $n$ на $\nu$:
\begin{equation}
Z_\nu^{(h)}=\sqrt{\frac{\varepsilon_{h+1}}{\varepsilon_h}}\,
\frac{C_\nu'+Z_\nu^{(h-1)}\,S_\nu'}{C_\nu+Z_\nu^{(h-1)}\,S_\nu},
\label{cone-cone-eq:adenkerker}
\end{equation}
где $C_\nu,S_\nu$ --- кросс-произведения функций Риккати--Бесселя нецелого индекса $\psi_\nu,\chi_\nu$ соседних слоёв. Радиальная рекурсия \emph{поканальна} по $\nu$ (разные $\nu$ внутри слоя не связаны, так как угловой базис фиксирован), и связь мод возникает лишь на финальной сшивке $r=a$ \eqref{cone-cone-eq:modematch-system}. Так стратификация добавляется без увеличения размерности связанной системы --- она лишь модифицирует эффективный импеданс конической области, подаваемый в \eqref{cone-cone-eq:modematch-system}.

\subsection{Путь к сферической $T$-матрице и сборка в GMM}
\label{cone-cone-ssec:tmatrix}

Результат сшивки \eqref{cone-cone-eq:modematch-system} --- блочная по $m$ матрица $\bm{T}^{(m)}$, действующая \emph{уже} в стандартном целочисленном сферическом VSWF-базисе внешней области (это и есть смысл проекции конических мод на $\{\bm M_{nm},\bm N_{nm}\}$ через $G^{(m)}_{n\nu}$). Собрав блоки по всем $m$, получаем полную $T$-матрицу тела
\begin{equation}
\bm{c}=\bm{T}^{\mathrm{cone}}\,\bm{d},
\qquad
T^{\mathrm{cone}}_{\sigma n m,\,\sigma' n' m'}=T^{(m)}_{\sigma n,\,\sigma' n'}\,\delta_{mm'},
\label{cone-cone-eq:Tcone}
\end{equation}
недиагональную по $n$ и смешивающую TE/TM ($\sigma,\sigma'\in\{M,N\}$) при $m\neq0$ --- в полном соответствии с общим определением $\bm c=\bm T\bm d$ \eqref{plan-eq:Tdef} и со структурой сфероидальной $T$-матрицы \eqref{spheroid-eq:Tsphere}. Поскольку \eqref{cone-cone-eq:Tcone} выражена в общем сферическом базисе, конус включается в суперпозицию $T$-матриц (GMM) точно так же, как сфера и сфероид: его $\bm{T}^{\mathrm{cone}}$ подставляется в самосогласованную систему \eqref{decomp-eq:gmm-system} наряду с трансляционными коэффициентами Штейна--Крузана \cite{stein1961,cruzan1962}, при условии непересечения описанных сфер \eqref{decomp-eq:non-overlap}.

\subsection{Дальнее поле и сечения}
\label{cone-cone-ssec:farfield}

Дальняя зона и сечения вычисляются из коэффициентов рассеяния $\bm c=\bm T^{\mathrm{cone}}\bm d$ в \emph{общем} сферическом базисе теми же формулами, что и для прочих примитивов \cite{mishchenko2002,bowman1987}. В $T$-матричной нормировке ($e^{-i\omega t}$, исходящая $h_n^{(1)}$) для фиксированной ориентации
\begin{equation}
C_{\mathrm{ext}}=-\frac{2\pi}{k_0^2}\,\mathrm{Re}\big[\bm{d}^{\dagger}\bm{T}^{\mathrm{cone}}\bm{d}\big],
\qquad
C_{\mathrm{sca}}=\frac{2\pi}{k_0^2}\,\big\|\bm{T}^{\mathrm{cone}}\bm{d}\big\|^2,
\qquad
C_{\mathrm{abs}}=C_{\mathrm{ext}}-C_{\mathrm{sca}},
\label{cone-cone-eq:cross}
\end{equation}
при единично нормированных амплитудах падающей волны; ориентационное усреднение даётся следами блоков $\bm{T}^{(m)}$ как в \eqref{spheroid-eq:cross} \cite{mishchenko2002}. Амплитудная диаграмма строится сверткой коэффициентов с целочисленными угловыми функциями $\pi_n,\tau_n$ внешней области; «конический» вклад вершины и ребра проявляется в медленно сходящихся компонентах ряда, стабилизируемых преобразованием последовательностей по образцу \cite{klinkenbusch2007}.

\subsection{Предельные переходы и проверки}
\label{cone-cone-ssec:limits}

Контроль корректности --- через геометрические пределы по полууглу $\theta_0$ и размеру $ka$.
\begin{itemize}
  \item \textbf{$\theta_0\to\pi$} (конус раскрывается почти до полного телесного угла): тело вырождается в сферический сегмент/шапку $r\le a$; собственные степени $\nu$ из \eqref{cone-cone-eq:eig-TM}--\eqref{cone-cone-eq:eig-TE} стремятся к целым $\nu\to n$ (полюс $\theta=0$ перестаёт быть исключённым), матрица перекрытия $G^{(m)}_{n\nu}\to\delta_{n\nu}$, связанная система \eqref{cone-cone-eq:modematch-system} становится блочно-диагональной по $n$, и $\bm{T}^{\mathrm{cone}}$ переходит в диагональную $T$-матрицу однородного шара (Ми) --- прямая редукция к каноническому ядру \S\ref{sphere-sec:sphere} \cite{bohren1983,mishchenko2002}.
  \item \textbf{$\theta_0\to0$} (тонкий конус/игла): доступный телесный угол стягивается, $\nu_{\min}\to\infty$, рассеяние стремится к рэлеевскому пределу тонкого вытянутого тела (доминирует осевая поляризуемость) --- согласуется с игольчатым пределом эллипсоида $L_a\to0$ \eqref{ellipsoid-eq:depol-sum} \cite{bohren1983}.
  \item \textbf{$ka\to0$}: квазистатика --- сечения должны воспроизводить дипольную поляризуемость конического тела (контроль через оптическую теорему $C_{\mathrm{abs}}=0$ для непоглощающего конуса).
  \item \textbf{Полу-бесконечный предел} ($a\to\infty$): сшивочные члены на $r=a$ исчезают, и решение сводится к точному коническому блоку \eqref{cone-cone-eq:greendyad} Клинкенбуша \cite{klinkenbusch2007} --- независимый аналитический эталон.
\end{itemize}

\subsection{Честная оценка: самый трудный примитив}
\label{cone-cone-ssec:honest}

Конус следует квалифицировать как полуаналитический примитив \textsc{TIER-B} с наибольшей реализационной стоимостью в библиотеке. Причины: (i)~\emph{нецелые присоединённые функции Лежандра} $P_\nu^m(\cos\theta)$ и поиск их собственных степеней $\nu$ из трансцендентных уравнений \eqref{cone-cone-eq:eig-TM}--\eqref{cone-cone-eq:eig-TE} требуют специальной численной машинерии (непрерывные дроби, гипергеометрические ряды), слабо или вовсе не покрытой стандартными библиотеками (в отличие от целочисленных $P_n^m$ ядра-сферы); (ii)~сшивка мод даёт \emph{связанную} линейную систему \eqref{cone-cone-eq:modematch-system}, недиагональную по $n$, со сходимостью, управляемой аспектным отношением и параметром размера $ka$; (iii)~\emph{две} сингулярные линии (вершина и ребро) требуют аккуратного условия Мейкснера и ускорения рядов \cite{klinkenbusch2007,bowman1987}. Это не «чистое» решение уровня сферы: ожидаются усечённая связанная система, контролируемая (а не машинная) точность и обязательная забота о ребре. Тем не менее путь строгий --- замкнутые собственные функции в естественных координатах, точный полу-бесконечный блок \cite{klinkenbusch2007} и стандартная сшивка/EBCM \cite{waterman1971} --- и итоговая $\bm{T}^{\mathrm{cone}}$ \eqref{cone-cone-eq:Tcone} полностью совместима с единым сферическим VSWF-каркасом GMM.

\section{Суперпозиция T-матриц и сложная геометрия}
\label{tmatrix-sec:tmatrix}

Аналитические решатели предыдущих разделов дают строгий отклик \emph{одиночного} конечного тела (сфера, сфероид, эллипсоид) в виде его T-матрицы в базисе сферических векторных волновых функций (ВВФ). Сложная геометрия в проекте GreenTensor собирается из этих примитивов методом суперпозиции T-матриц --- обобщённым методом Ми для совокупности частиц (Generalized Multiparticle Mie, GMM) \cite{xu1995,mishchenko2002}. Метод строг в рамках полного волнового уравнения Максвелла: он не вводит приближений сверх усечения мультипольных рядов и сходится точно при выполнении геометрического условия раздельности описывающих сфер (см.\ \S\ref{tmatrix-sec:tm-conv}).

\subsection{Регулярные и исходящие векторные сферические волновые функции}
\label{tmatrix-sec:tm-vswf}

Зафиксируем временну\'ю конвенцию $e^{-i\omega t}$ и волновое число фоновой (внешней) среды $k$. Скалярные мультипольные решения уравнения Гельмгольца обозначим
\begin{equation}
\psi_{mn}^{(J)}(\bm{r}) = z_n^{(J)}(kr)\, P_n^{|m|}(\cos\theta)\, e^{i m \varphi},
\qquad n\ge 1,\ |m|\le n,
\label{tmatrix-eq:tm-scalar}
\end{equation}
где $z_n^{(1)}=j_n$ (регулярная радиальная зависимость, индекс $J{=}1$) и $z_n^{(3)}=h_n^{(1)}$ (исходящая волна, индекс $J{=}3$), согласованная с конвенцией $e^{-i\omega t}$ \cite{stratton1941,bohren1983}. Векторные сферические волновые функции определяются стандартно \cite{stratton1941,mishchenko2002}:
\begin{align}
\bm{M}_{mn}^{(J)}(\bm{r}) &= \nabla\times\bigl[\bm{r}\,\psi_{mn}^{(J)}(\bm{r})\bigr],
\label{tmatrix-eq:tm-M}\\
\bm{N}_{mn}^{(J)}(\bm{r}) &= \frac{1}{k}\,\nabla\times\bm{M}_{mn}^{(J)}(\bm{r}).
\label{tmatrix-eq:tm-N}
\end{align}
Регулярные функции (конечные в начале координат) обозначим $\mathrm{Rg}\,\bm{M}_{mn}=\bm{M}_{mn}^{(1)}$, $\mathrm{Rg}\,\bm{N}_{mn}=\bm{N}_{mn}^{(1)}$; исходящие (удовлетворяющие условию излучения Зоммерфельда) --- $\bm{M}_{mn}=\bm{M}_{mn}^{(3)}$, $\bm{N}_{mn}=\bm{N}_{mn}^{(3)}$. Здесь $\bm{M}$ --- TE-(магнитный) тип, $\bm{N}$ --- TM-(электрический) тип относительно радиального направления; именно эти два типа отвечают за коэффициенты $M_n$ (TM, импеданс $Z$) и $N_n$ (TE, адмиттанс $Y$) диагональной матрицы Ми из ядра \texttt{mie\_core} (раздел о сфере).

\subsection{T-матрица одиночного тела}
\label{tmatrix-sec:tm-single}

Падающее на тело поле раскладывается по \emph{регулярным} ВВФ, рассеянное --- по \emph{исходящим}:
\begin{align}
\bm{E}^{\mathrm{inc}}(\bm{r}) &= \sum_{n=1}^{\infty}\sum_{m=-n}^{n}
\Bigl[\, d_{mn}^{(1)}\,\mathrm{Rg}\,\bm{M}_{mn}(\bm{r}) + d_{mn}^{(2)}\,\mathrm{Rg}\,\bm{N}_{mn}(\bm{r}) \Bigr],
\label{tmatrix-eq:tm-Einc}\\
\bm{E}^{\mathrm{sca}}(\bm{r}) &= \sum_{n=1}^{\infty}\sum_{m=-n}^{n}
\Bigl[\, c_{mn}^{(1)}\,\bm{M}_{mn}(\bm{r}) + c_{mn}^{(2)}\,\bm{N}_{mn}(\bm{r}) \Bigr].
\label{tmatrix-eq:tm-Esca}
\end{align}
Собрав коэффициенты в столбцы $\bm{c}=(c_{mn}^{(1)},c_{mn}^{(2)})^{\!\top}$ и $\bm{d}=(d_{mn}^{(1)},d_{mn}^{(2)})^{\!\top}$, линейность задачи рассеяния даёт связь через T-матрицу \cite{waterman1971,mishchenko2002}:
\begin{equation}
\bm{c} = \mathbf{T}\,\bm{d},
\qquad
\mathbf{T}=\begin{pmatrix} \mathbf{T}^{11} & \mathbf{T}^{12}\\[2pt] \mathbf{T}^{21} & \mathbf{T}^{22}\end{pmatrix}.
\label{tmatrix-eq:tm-Td}
\end{equation}
Для сферы $\mathbf{T}$ диагональна и не смешивает TE/TM: $T_{mn,m'n'}^{11}=-N_n\,\delta_{nn'}\delta_{mm'}$, $T_{mn,m'n'}^{22}=-M_n\,\delta_{nn'}\delta_{mm'}$ (знак и нормировка фиксируются конвенцией коэффициентов Ми из \texttt{mie\_core}); это и есть классический результат Ми и Адена--Керкера \cite{aden1951,bohren1983,mishchenko2002}. Для конфокального многослойного сфероида $\mathbf{T}$ блочно-диагональна по $m$, но связывает порядки $n$ и типы M/N; её получают пересчётом сфероидального решения в сферический базис (см.\ раздел о сфероиде). В любом случае на вход GMM каждое тело $i$ поступает как готовая матрица $\mathbf{T}_i$ в его собственной (центрированной и ориентированной) системе.

\subsection{Теорема сложения (трансляция) для векторных волновых функций}
\label{tmatrix-sec:tm-addition}

Ключевой ингредиент --- пересчёт разложения, центрированного в точке $O_j$, в разложение относительно другой точки $O_i$. Пусть $\bm{r}_{ij}=\bm{r}_i-\bm{r}_j$ соединяет центры, а наблюдательный радиус-вектор связан соотношением $\bm{r}_j=\bm{r}_i+\bm{r}_{ij}$. Теорема трансляционного сложения ВВФ \cite{stein1961,cruzan1962,mishchenko2002} выражает функции в системе $j$ через функции в системе $i$:
\begin{align}
\bm{M}_{\mu\nu}^{(p)}(\bm{r}_j)
&= \sum_{n,m}\Bigl[\, A_{\mu\nu}^{mn}(\bm{r}_{ij})\,\bm{M}_{mn}^{(q)}(\bm{r}_i)
+ B_{\mu\nu}^{mn}(\bm{r}_{ij})\,\bm{N}_{mn}^{(q)}(\bm{r}_i)\Bigr],
\label{tmatrix-eq:tm-addM}\\
\bm{N}_{\mu\nu}^{(p)}(\bm{r}_j)
&= \sum_{n,m}\Bigl[\, A_{\mu\nu}^{mn}(\bm{r}_{ij})\,\bm{N}_{mn}^{(q)}(\bm{r}_i)
+ B_{\mu\nu}^{mn}(\bm{r}_{ij})\,\bm{M}_{mn}^{(q)}(\bm{r}_i)\Bigr].
\label{tmatrix-eq:tm-addN}
\end{align}
Структура коэффициентов одинакова для $\bm{M}$ и $\bm{N}$, с перестановкой ролей $A\leftrightarrow B$: коэффициент $A_{\mu\nu}^{mn}$ связывает функции одного типа (M$\to$M, N$\to$N), $B_{\mu\nu}^{mn}$ --- разных (M$\to$N, N$\to$M); это в точности структура (C.66)--(C.70) и (5.229) у Мищенко \cite{mishchenko2002}. Выбор радиальной зависимости результата задаётся индексом $q$ относительно исходного $p$:
\begin{itemize}
\item \textbf{регулярная трансляция} ($p=q=1$, $\mathrm{Rg}\to\mathrm{Rg}$): коэффициенты $A,B$ строятся на сферических функциях Бесселя $j_l(k r_{ij})$; используется для пересчёта \emph{падающего} поля между центрами;
\item \textbf{исходящая$\to$регулярная} трансляция ($p=3$, $q=1$): коэффициенты $A,B$ строятся на $h_l^{(1)}(k r_{ij})$; используется для пересчёта поля, \emph{рассеянного} телом $j$ (исходящего вокруг $O_j$), в \emph{падающее} (регулярное в окрестности $O_i$, при $r_i<r_{ij}$) на тело $i$.
\end{itemize}
В замкнутой форме Крузана \cite{cruzan1962} коэффициенты выражаются через коэффициенты Гонта (интегралы трёх присоединённых функций Лежандра) и сферические гармоники направления $\hat{\bm{r}}_{ij}$:
\begin{align}
A_{\mu\nu}^{mn} &= (-1)^{m}\,\frac{i^{\,\nu-n}\,(2n+1)}
{2\,n(n+1)}\,
\frac{(n+m)!\,(\nu-\mu)!}{(n-m)!\,(\nu+\mu)!}
\sum_{l}\, i^{\,l}\,
\bigl[\,\nu(\nu+1)+n(n+1)-l(l+1)\bigr]\,
a(m,n,-\mu,\nu,l)\;
z_l^{(p)}(k r_{ij})\,Y_{l}^{\,m-\mu}(\hat{\bm{r}}_{ij}),
\label{tmatrix-eq:tm-Acoef}\\[4pt]
B_{\mu\nu}^{mn} &= (-1)^{m+1}\,\frac{i^{\,\nu-n}\,(2n+1)}
{2\,n(n+1)}\,
\frac{(n+m)!\,(\nu-\mu)!}{(n-m)!\,(\nu+\mu)!}
\sum_{l}\, i^{\,l}\,
b(m,n,-\mu,\nu,l)\,
z_l^{(p)}(k r_{ij})\,Y_{l}^{\,m-\mu}(\hat{\bm{r}}_{ij}),
\label{tmatrix-eq:tm-Bcoef}
\end{align}
где $z_l^{(1)}=j_l$, $z_l^{(3)}=h_l^{(1)}$; $a(\cdot)$ --- линеаризационные (Гонта) коэффициенты, а $b(\cdot)$ --- их «нечётная» по чётности модификация. Суммирование по $l$ идёт по правилам отбора Гонта: $|n-\nu|\le l\le n+\nu$, с шагом $2$ для $a$ и сдвигом на $1$ для $b$, и $l\ge |m-\mu|$. Конкретный вид нормировки $a,b$ и эквивалентная запись Штейна \cite{stein1961} приведены в первоисточниках; для реализации удобны рекуррентные соотношения Сюя \cite{xu1995}, дающие $A_{\mu\nu}^{mn},B_{\mu\nu}^{mn}$ численно устойчиво без прямого вычисления факториалов. Введём блочный \emph{трансляционный оператор}
\begin{equation}
\bm{\beta}(\bm{r}_{ij})=
\begin{pmatrix} \mathbf{A}(\bm{r}_{ij}) & \mathbf{B}(\bm{r}_{ij})\\[2pt] \mathbf{B}(\bm{r}_{ij}) & \mathbf{A}(\bm{r}_{ij})\end{pmatrix},
\label{tmatrix-eq:tm-beta}
\end{equation}
действующий на сдвоенные столбцы $(c^{(1)},c^{(2)})$ так же, как \eqref{tmatrix-eq:tm-addM}--\eqref{tmatrix-eq:tm-addN}; это матрица переноса из (5.229)--(5.230) у Мищенко \cite{mishchenko2002}. Обозначим $\bm{\beta}^{(3)}$ оператор с $z_l^{(3)}=h_l^{(1)}$ (исходящий$\to$регулярный сдвиг) и $\bm{\beta}^{(1)}$ --- регулярный ($z_l^{(1)}=j_l$).

\subsection{Явная замкнутая форма коэффициентов \texorpdfstring{$A,B$}{A,B} (реализованная)}
\label{tmatrix-sec:tm-explicit}

Абстрактные линеаризационные коэффициенты $a(\cdot),b(\cdot)$ в \eqref{tmatrix-eq:tm-Acoef}--\eqref{tmatrix-eq:tm-Bcoef} ниже конкретизированы до явных замкнутых выражений, реализованных в библиотеке (\texttt{vswf.translation\_block\_closed}). Введём общий радиально-угловой множитель (тот же «скелет», что и в скалярной теореме сложения)
\begin{equation}
\tau_l \equiv 4\pi\, i^{\,n-\nu}\,(-1)^m\,(-i)^{\,l}\, z_l^{(p)}(k r_{ij})\,
\bigl[\,Y_l^{\,m-\mu}(\hat{\bm r}_{ij})\,\bigr]^{*},
\qquad z_l^{(1)}=j_l,\ \ z_l^{(3)}=h_l^{(1)},
\label{tmatrix-eq:tm-tau}
\end{equation}
и коэффициент Гонта в нормировке через символы Вигнера $3j$:
\begin{equation}
\mathcal{G}^{\,m\mu}_{n\nu l}=\sqrt{\frac{(2n+1)(2\nu+1)(2l+1)}{4\pi}}\,
\begin{pmatrix} n&\nu&l\\ 0&0&0\end{pmatrix}
\begin{pmatrix} n&\nu&l\\ -m&\mu&m-\mu\end{pmatrix}.
\label{tmatrix-eq:tm-gaunt}
\end{equation}
Тогда (источник $(\nu,\mu)$, цель $(n,m)$; конвенция $e^{-i\omega t}$, исходящая волна $\sim h_l^{(1)}$):
\begin{align}
A_{\mu\nu}^{mn} &= \frac{1}{2\,n(n+1)}
\!\!\sum_{\substack{l=|n-\nu|\\[1pt] n+\nu+l\ \mathrm{чётн.}}}^{\,n+\nu}\!\!
\bigl[\,\nu(\nu+1)+n(n+1)-l(l+1)\,\bigr]\,
\mathcal{G}^{\,m\mu}_{n\nu l}\,\tau_l,
\label{tmatrix-eq:tm-Aexplicit}\\[4pt]
B_{\mu\nu}^{mn} &=
\!\!\sum_{\substack{l=|n-\nu|\\[1pt] n+\nu+l\ \mathrm{нечётн.}}}^{\,n+\nu}\!\!
W_{n\nu l}\,
\begin{pmatrix} n&\nu&l\\ -m&\mu&m-\mu\end{pmatrix}
\begin{pmatrix} n&\nu&l-1\\ 0&0&0\end{pmatrix}\tau_l,
\label{tmatrix-eq:tm-Bexplicit}
\end{align}
где аналитический вес коэффициента $B$ есть
\begin{equation}
W_{n\nu l}=-\,\frac{\sqrt{\,(2n+1)(2\nu+1)(2l+1)\,(n{+}\nu{+}l{+}1)(n{+}\nu{-}l{+}1)(l{+}n{-}\nu)(l{-}n{+}\nu)\,}}
{4\sqrt{\pi}\,\,n(n+1)}.
\label{tmatrix-eq:tm-Wbeta}
\end{equation}

\paragraph{Происхождение и проверка.} Весовой множитель $A$ в \eqref{tmatrix-eq:tm-Aexplicit}, $[\nu(\nu{+}1)+n(n{+}1)-l(l{+}1)]/[2n(n{+}1)]$, согласуется с операторной формой Виттмана \cite{wittmann1988} (для $n=\nu$ --- тождественно; нормировка целевой моды $1/[2n(n{+}1)]$ зафиксирована сверкой). Аналитический вес $B$ \eqref{tmatrix-eq:tm-Wbeta} получен распознаванием из точной таблицы коэффициентов: после выделения универсального множителя $(2n{+}1)(2\nu{+}1)/[3(2l{+}1)]$ выражение сворачивается к симметричной по $(n,\nu)$ замкнутой форме. Обе формулы верифицированы независимо: (i) против спектрального вычисления тех же коэффициентов проекцией на сферу --- совпадение $\sim\!10^{-9}$ для $A$ и $\sim\!10^{-11}$ для $B$; (ii) как тождество полей $\bm M^{(3)}_{\mu\nu}(\bm r-\bm r_{ij})=\sum_{nm}[\,A_{\mu\nu}^{mn}\,\mathrm{Rg}\,\bm M_{mn}+B_{\mu\nu}^{mn}\,\mathrm{Rg}\,\bm N_{mn}\,]$ (и аналогично для $\bm N$) в контрольных точках при \emph{произвольном} разносе --- невязка $\sim\!10^{-5}$ при $n_{\max}=12$, монотонно убывающая с $n_{\max}$. В отличие от спектральной проекции, замкнутые \eqref{tmatrix-eq:tm-Aexplicit}--\eqref{tmatrix-eq:tm-Bexplicit} \emph{не} имеют численного ограничения $k r_{ij}\gtrsim n_{\max}$ и точны всюду, где сходится ряд \eqref{tmatrix-eq:tm-addM}--\eqref{tmatrix-eq:tm-addN} ($r_i<r_{ij}$); это снимает потерю точности на плотных кластерах (касающиеся примитивы при разложении сложной геометрии).

\subsection{Самосогласованная система GMM}
\label{tmatrix-sec:tm-system}

Пусть имеется $L$ тел с центрами $\bm{r}_i$ и T-матрицами $\mathbf{T}_i$. Поле, \emph{возбуждающее} тело $i$, есть сумма внешнего падающего поля и полей, \emph{рассеянных всеми остальными} телами и пересчитанных к центру $O_i$. Поскольку рассеянные поля исходящие, их перенос в регулярное разложение вокруг $O_i$ выполняется оператором $\bm{\beta}^{(3)}(\bm{r}_{ij})$. Обозначим $\bm{d}_i$ --- коэффициенты падающего (внешнего) поля, разложенного вокруг $O_i$, $\bm{a}_i$ --- коэффициенты полного возбуждающего поля, $\bm{c}_i=\mathbf{T}_i\bm{a}_i$ --- рассеянного. Тогда \cite{xu1995,mishchenko2002}
\begin{equation}
\bm{a}_i \;=\; \bm{d}_i \;+\; \sum_{j\ne i}^{L}\bm{\beta}^{(3)}(\bm{r}_{ij})\,\bm{c}_j,
\qquad
\bm{c}_i \;=\; \mathbf{T}_i\,\bm{a}_i,
\qquad i=1,\dots,L.
\label{tmatrix-eq:tm-selfcons}
\end{equation}
Подстановка $\bm{c}_j=\mathbf{T}_j\bm{a}_j$ даёт замкнутую систему относительно возбуждающих коэффициентов:
\begin{equation}
\boxed{\;\bm{a}_i - \sum_{j\ne i}\bm{\beta}^{(3)}(\bm{r}_{ij})\,\mathbf{T}_j\,\bm{a}_j \;=\; \bm{d}_i\;}
\qquad i=1,\dots,L,
\label{tmatrix-eq:tm-linsys}
\end{equation}
или в блочно-матричной форме для составного вектора $\bm{a}=(\bm{a}_1,\dots,\bm{a}_L)^{\!\top}$:
\begin{equation}
\bigl(\mathbf{I}-\mathbf{H}\bigr)\,\bm{a}=\bm{d},
\qquad
\mathbf{H}_{ij}=
\begin{cases}
\bm{\beta}^{(3)}(\bm{r}_{ij})\,\mathbf{T}_j, & i\ne j,\\[2pt]
\mathbf{0}, & i=j.
\end{cases}
\label{tmatrix-eq:tm-block}
\end{equation}
Это в точности система суперпозиции T-матриц (5.230) у Мищенко \cite{mishchenko2002}. Формальное решение
\begin{equation}
\bm{a}=(\mathbf{I}-\mathbf{H})^{-1}\bm{d},
\qquad
\bm{c}_i=\mathbf{T}_i\,\bm{a}_i,
\label{tmatrix-eq:tm-solve}
\end{equation}
эквивалентно борновскому ряду многократного рассеяния $\bm{a}=\sum_{p\ge0}\mathbf{H}^p\bm{d}$: член $p$ описывает $p$-кратное переотражение между телами. На практике \eqref{tmatrix-eq:tm-block} решают прямым обращением (LU) для умеренного $L$ либо итерационно (порядок Неймана/GMRES) при больших кластерах \cite{xu1995}. Усечение по $n\le n_{\max,i}$ для каждого тела определяется его параметром размера и расстояниями (правило типа Уискомба плюс запас на сходимость трансляций).

Из решённых $\bm{c}_i$ можно построить \emph{составную} (кластерную) T-матрицу относительно единого начала $O$ \cite{mishchenko2002}: перенося падающее поле к центрам тел оператором $\bm{\beta}^{(1)}(\bm{r}_{jO})$ и собирая рассеянное обратно к $O$ регулярным оператором $\bm{\beta}^{(1)}(\bm{r}_{Oi})$ (формула (5.236) у Мищенко),
\begin{equation}
\mathbf{T}^{\mathrm{cluster}}
=\sum_{i=1}^{L}\sum_{j=1}^{L}
\bm{\beta}^{(1)}(\bm{r}_{Oi})\,\mathbf{T}_{ij}\,\bm{\beta}^{(1)}(\bm{r}_{jO}),
\qquad
\mathbf{T}_{ij}=\mathbf{T}_i\bigl[(\mathbf{I}-\mathbf{H})^{-1}\bigr]_{ij},
\label{tmatrix-eq:tm-cluster}
\end{equation}
которая ориентационно-независима и допускает дальнейшее использование в качестве единого «примитива» более высокого уровня.

\subsection{Дальнее поле и сечения}
\label{tmatrix-sec:tm-farfield}

В дальней зоне ($kr\to\infty$) исходящие ВВФ дают $\bm{M}_{mn}^{(3)},\bm{N}_{mn}^{(3)}\sim (-i)^{n}\,\dfrac{e^{ikr}}{kr}\,\bm{\mathcal{X}}_{mn}(\hat{\bm{r}})$ с известными угловыми векторами $\bm{\mathcal{X}}_{mn}$ (поперечные гармоники $\bm{B}_{mn},\bm{C}_{mn}$ в \cite{mishchenko2002}, ср.\ (5.9)); множитель $i^{-n}=(-i)^n$ отвечает за асимптотику $h_n^{(1)}(kr)\sim(-i)^{n+1}e^{ikr}/(kr)$ \cite{bohren1983,mishchenko2002,dlmf}. Полная амплитуда рассеяния кластера складывается из вкладов всех тел с \emph{фазовыми множителями}, учитывающими смещение центров $\bm{r}_i$ относительно общего начала:
\begin{equation}
\bm{E}^{\mathrm{sca}}(\bm{r})\xrightarrow[kr\to\infty]{}
\frac{e^{ikr}}{kr}\,\bm{F}(\hat{\bm{r}}),\qquad
\bm{F}(\hat{\bm{r}})=\sum_{i=1}^{L} e^{-ik\,\hat{\bm{r}}\cdot\bm{r}_i}
\sum_{n,m}(-i)^{n}\Bigl[c_{i,mn}^{(1)}\,\bm{\mathcal{X}}^{M}_{mn}(\hat{\bm{r}})+c_{i,mn}^{(2)}\,\bm{\mathcal{X}}^{N}_{mn}(\hat{\bm{r}})\Bigr].
\label{tmatrix-eq:tm-farfield}
\end{equation}
Множитель $e^{-ik\,\hat{\bm{r}}\cdot\bm{r}_i}$ --- фаза, набегающая из-за положения тела $i$ (раскрытие $k|\bm{r}-\bm{r}_i|\approx kr-k\,\hat{\bm{r}}\cdot\bm{r}_i$); именно интерференция этих фаз формирует структуру диаграммы рассеяния кластера (бистатическая RCS). Сечение рассеяния для падающей плоской волны единичной амплитуды получается интегрированием квадрата амплитуды по полному телесному углу, что для ортонормированных угловых функций сводится к сумме квадратов коэффициентов \cite{mishchenko2002}:
\begin{equation}
C_{\mathrm{sca}} = \frac{1}{k^2}\oint |\bm{F}(\hat{\bm{r}})|^2\,d\Omega
= \frac{1}{k^2}\sum_{i,j}\sum_{\substack{n,m\\ \nu,\mu}}
\Bigl[\,c_{i,mn}^{(1)}\bigl(c_{j,\mu\nu}^{(1)}\bigr)^{\!*}\,I^{MM}_{mn,\mu\nu}
+ c_{i,mn}^{(2)}\bigl(c_{j,\mu\nu}^{(2)}\bigr)^{\!*}\,I^{NN}_{mn,\mu\nu}+\dots\Bigr],
\label{tmatrix-eq:tm-Csca}
\end{equation}
где $I_{\cdot}$ --- угловые интегралы перекрытия гармоник, учитывающие межчастичные фазы $e^{-ik\hat{\bm{r}}\cdot(\bm{r}_i-\bm{r}_j)}$ (для одиночного тела $I^{MM}_{mn,\mu\nu}=I^{NN}_{mn,\mu\nu}\propto\delta_{nm,\nu\mu}$ и \eqref{tmatrix-eq:tm-Csca} переходит в $C_{\mathrm{sca}}=k^{-2}\sum|c_{mn}|^2$, ср.\ (5.18b)). Сечение экстинкции даётся оптической теоремой через прямое (forward, $\hat{\bm{r}}=\hat{\bm{k}}_{\mathrm{inc}}$) рассеяние. В коэффициентной форме (5.18a) Мищенко \cite{mishchenko2002}, с конвенцией $e^{-i\omega t}$ и нормировкой коэффициентов \eqref{tmatrix-eq:tm-Einc}--\eqref{tmatrix-eq:tm-Esca},
\begin{equation}
C_{\mathrm{ext}} = -\frac{1}{k^2}\,\mathrm{Re}\sum_{i=1}^{L}\sum_{n,m}
\Bigl[\, d_{i,mn}^{(1)}\bigl(c_{i,mn}^{(1)}\bigr)^{\!*}
+ d_{i,mn}^{(2)}\bigl(c_{i,mn}^{(2)}\bigr)^{\!*}\Bigr]
\;=\;\frac{4\pi}{k^2}\,\mathrm{Re}\!\left[\hat{\bm{e}}^{*}_{\mathrm{inc}}\cdot\bm{F}(\hat{\bm{k}}_{\mathrm{inc}})\right],
\label{tmatrix-eq:tm-Cext}
\end{equation}
где $d_{i,mn}$ --- коэффициенты падающего поля, разложенного вокруг $O_i$; знак минус в коэффициентной форме фиксируется конвенцией $e^{-i\omega t}$ (формула (5.18a) у Мищенко), а эквивалентная запись через прямую амплитуду $\bm{F}$ есть классическая оптическая теорема \cite{bohren1983,mishchenko2002}. Сечение поглощения --- $C_{\mathrm{abs}}=C_{\mathrm{ext}}-C_{\mathrm{sca}}\ge 0$ \cite{bohren1983,mishchenko2002}. Эти выражения сводятся к диагональным формулам Ми из \texttt{mie\_core} ($C_{\mathrm{ext}}=\tfrac{2\pi}{k^2}\sum_n(2n+1)\,\mathrm{Re}(\dots)$, $C_{\mathrm{sca}}=\tfrac{2\pi}{k^2}\sum_n(2n+1)(\dots)$, ср.\ (5.156)--(5.157) у Мищенко) при $L=1$ и сферической $\mathbf{T}$, что служит проверкой согласованности.

\subsection{Ориентация тел: вращения Вигнера}
\label{tmatrix-sec:tm-rotation}

T-матрица $\mathbf{T}_i^{\mathrm{body}}$ вычисляется в собственной (главной) системе тела. Если тело повёрнуто на углы Эйлера $(\alpha_i,\beta_i,\gamma_i)$ относительно лабораторной системы кластера, её приводят к лабораторным осям подобием с матрицами вращения ВВФ \cite{mishchenko2002}:
\begin{equation}
\mathbf{T}_i^{\mathrm{lab}} = \mathbf{D}(\alpha_i,\beta_i,\gamma_i)\,\mathbf{T}_i^{\mathrm{body}}\,\mathbf{D}^{-1}(\alpha_i,\beta_i,\gamma_i),
\label{tmatrix-eq:tm-rot}
\end{equation}
где блоки $\mathbf{D}$ диагональны по $n$ и составлены из элементов матрицы Вигнера
\begin{equation}
D_{m'm}^{n}(\alpha,\beta,\gamma)=e^{-im'\alpha}\,d_{m'm}^{n}(\beta)\,e^{-im\gamma},
\label{tmatrix-eq:tm-wignerD}
\end{equation}
$d_{m'm}^n(\beta)$ --- малая (вещественная) d-функция Вигнера. Поскольку вращение сохраняет $n$ и не смешивает M/N-типы (это свойство ВВФ при поворотах), \eqref{tmatrix-eq:tm-rot} эффективно и применяется до сборки $\mathbf{H}$. Для сферы $\mathbf{D}\mathbf{T}\mathbf{D}^{-1}=\mathbf{T}$ (изотропия), поэтому вращения существенны лишь для сфероидов/эллипсоидов и для произвольно ориентированных составных примитивов.

\subsection{Условие сходимости}
\label{tmatrix-sec:tm-conv}

Перенос рассеянного поля тела $j$ в регулярное разложение вокруг $O_i$ оператором $\bm{\beta}^{(3)}(\bm{r}_{ij})$ математически строг \emph{лишь там, где это разложение сходится} --- то есть внутри наибольшей сферы с центром $O_i$, не содержащей особенностей поля тела $j$ (условие $r_i<r_{ij}$ в \eqref{tmatrix-eq:tm-addM}--\eqref{tmatrix-eq:tm-addN}) \cite{stein1961,cruzan1962,mishchenko2002}. Отсюда фундаментальное условие применимости GMM:

\begin{quote}
\textbf{Метод суперпозиции T-матриц строг тогда и только тогда, когда описывающие сферы тел не пересекаются (попарно непересекающиеся circumscribing spheres).}
\end{quote}

Формально, для каждой пары $i\ne j$ должно выполняться $r_{ij} > R_i^{\mathrm{circ}} + R_j^{\mathrm{circ}}$ --- расстояние между центрами больше суммы радиусов описывающих сфер. При выполнении этого условия ряд многократного рассеяния \eqref{tmatrix-eq:tm-solve} сходится, а результат точен в пределе $n_{\max}\to\infty$ \cite{mishchenko2002}. Практические следствия:
\begin{itemize}
\item \textbf{Непересекающиеся тела} --- решение строгое; скорость сходимости по $n_{\max}$ тем выше, чем больше зазор между описывающими сферами.
\item \textbf{Касающиеся тела} ($r_{ij}\approx R_i^{\mathrm{circ}}+R_j^{\mathrm{circ}}$) --- ряд сходится медленно, точка касания лежит на границе области сходимости; требуется большой $n_{\max}$ и контроль обусловленности $\mathbf{I}-\mathbf{H}$.
\item \textbf{Пересекающиеся описывающие сферы} (даже если сами тела геометрически не перекрываются, например два сильно вытянутых сфероида «бок о бок») --- классический GMM в форме \eqref{tmatrix-eq:tm-block} теряет строгую сходимость; необходимы расширения (промежуточные центры, метод нескольких разложений, либо гибрид с поверхностными интегральными уравнениями).
\end{itemize}
Поэтому в проекте сложная геометрия конструируется из примитивов, чьи описывающие сферы раздельны; для конфигураций с близким касанием контроль точности ведётся по невязке системы \eqref{tmatrix-eq:tm-block} и по выполнению оптической теоремы $C_{\mathrm{abs}}\ge 0$ как физического индикатора сходимости \cite{mishchenko2002}.

\section{Разложение сложной геометрии и протокол верификации}
\label{decomp-sec:decomp}

Настоящий раздел замыкает теоретическую конструкцию проекта. В первой части описано, каким образом тело произвольной (в том числе «угловатой», коробчатой) формы редуцируется к \emph{кластеру аналитических примитивов} --- слоистых сфер, конфокальных сфероидов и (квазистатических) эллипсоидов, --- и каким образом такое разложение является точным, а где носит характер управляемой аппроксимации. Вторая часть формулирует строгий протокол верификации: каждый решатель сопоставляется с независимым аналитическим пределом и с опубликованным эталоном, а сквозными (не зависящими от конкретного решателя) критериями выступают закон сохранения энергии (оптическая теорема) и взаимность.

\subsection{Разложение сложной геометрии на аналитические примитивы}
\label{decomp-ssec:decomp-geom}

\paragraph{Принцип суперпозиции T-матриц (GMM).}
Пусть рассеиватель представлен как объединение $P$ непересекающихся примитивов с центрами $\bm{r}_p$, $p=1,\dots,P$. Полное поле есть сумма падающего и рассеянных каждым примитивом полей,
\begin{equation}
\bm{E} = \bm{E}^{\mathrm{inc}} + \sum_{p=1}^{P}\bm{E}^{\mathrm{sca}}_p ,
\label{decomp-eq:superpos}
\end{equation}
причём «возбуждение», действующее на $p$-й примитив, складывается из внешнего падающего поля и полей, рассеянных всеми остальными примитивами. В базисе векторных сферических волновых функций (VSWF) каждый примитив характеризуется своей T-матрицей $\mathsf{T}^{(p)}$, связывающей коэффициенты возбуждения $\bm{a}^{(p)}$ и рассеяния $\bm{f}^{(p)}=\mathsf{T}^{(p)}\bm{a}^{(p)}$. Самосогласованная система метода обобщённой мультисферы (Generalized Multiparticle Mie, GMM) имеет вид
\begin{equation}
\bm{f}^{(p)} = \mathsf{T}^{(p)}\Bigl(\bm{a}^{(p)}_{\mathrm{inc}} + \sum_{q\neq p}\mathsf{O}^{(pq)}\,\bm{f}^{(q)}\Bigr),
\qquad p=1,\dots,P,
\label{decomp-eq:gmm-system}
\end{equation}
где $\mathsf{O}^{(pq)}$ --- матрица трансляционного сложения (regular-to-irregular re-expansion), переносящая исходящие сферические гармоники из центра $q$ в регулярные гармоники с центром $p$ \cite{xu1995,stein1961,cruzan1962}. Коэффициенты $\mathsf{O}^{(pq)}$ суть трансляционные коэффициенты сложения сферических векторных волн (теоремы Штейна и Крузана) \cite{stein1961,cruzan1962,mishchenko2002}. Система \eqref{decomp-eq:gmm-system} решается прямо или итеративно (по порядку многократного рассеяния); сходимость гарантирована при невырожденности оператора $\mathsf{I}-\mathsf{T}\mathsf{O}$.

\paragraph{T-матрицы примитивов в общем базисе.}
Каждый примитив поставляет свою $\mathsf{T}^{(p)}$ в \emph{едином} сферическом VSWF-базисе:
\begin{itemize}
  \item \textbf{слоистая сфера} --- $\mathsf{T}^{(p)}$ диагональна, $T^{m}_{n}=\{-b_n,-a_n\}$ (коэффициенты Ми $a_n,b_n$ из импедансно-адмитансной рекурсии Адена--Керкера) \cite{aden1951,bohren1983};
  \item \textbf{конфокальный слоистый сфероид} --- $\mathsf{T}^{(p)}$ строится в собственном сфероидальном базисе (блочная по азимутальному $m$ рекурсия) и \emph{перепроецируется} на сферический базис для подстановки в \eqref{decomp-eq:gmm-system} \cite{asano1975,farafonov2012};
  \item \textbf{триаксиальный эллипсоид} --- аналитически доступен лишь квазистатический предел: примитив входит в \eqref{decomp-eq:gmm-system} как точечный диполь с тензором поляризуемости $\bm{\alpha}$, определяемым факторами деполяризации $L_x,L_y,L_z$ \cite{sihvola1990,bohren1983}.
\end{itemize}

\paragraph{Условие непересечения описанных сфер.}
Трансляционные разложения $\mathsf{O}^{(pq)}$ сходятся, если регулярное и иррегулярное разложения определены в непересекающихся областях; на практике это требование \emph{непересечения описанных сфер}: для любой пары примитивов
\begin{equation}
\lVert \bm{r}_p - \bm{r}_q\rVert \;>\; R^{\mathrm{circ}}_p + R^{\mathrm{circ}}_q,
\label{decomp-eq:non-overlap}
\end{equation}
где $R^{\mathrm{circ}}_p$ --- радиус наименьшей сферы, описанной вокруг $p$-го примитива \cite{xu1995,mishchenko2002}. При выполнении \eqref{decomp-eq:non-overlap} разложение по примитивам и решение \eqref{decomp-eq:gmm-system} являются \emph{строго точными} (в пределах усечения мультипольного ряда $n\le n_{\max}$): погрешность контролируется только числом удержанных гармоник и не содержит геометрической аппроксимации формы.

\paragraph{Что точно, а что приближённо.}
Разделим источники погрешности:
\begin{itemize}
  \item \emph{Точно} (до усечения ряда): отдельный примитив из разрешённого семейства; кластер примитивов, удовлетворяющий \eqref{decomp-eq:non-overlap}; слоистость сферы (произвольные радиусы) и конфокальная слоистость сфероида.
  \item \emph{Приближённо}: (i) аппроксимация формы реального тела набором примитивов --- геометрическая невязка границы; (ii) случай близкого контакта или взаимного перекрытия описанных сфер \eqref{decomp-eq:non-overlap}, когда трансляционные ряды теряют сходимость и требуется специальная регуляризация; (iii) электрически крупный или сильноконтрастный эллипсоид, для которого правомерен лишь квазистатический (релеевский) предел \cite{sihvola1990}.
\end{itemize}

\paragraph{Коробчатые («блочные») тела.}
Принципиальная установка проекта --- \emph{не} вводить отдельный «решатель параллелепипеда». Угловатое тело аппроксимируется композицией разрешённых примитивов (например, сильно сплюснутый/вытянутый сфероид как несущий блок плюс окаймляющие сферы), а острые рёбра и углы --- областью, где сферический мультипольный ряд сходится медленнее и требует увеличения $n_{\max}$. Такой подход сохраняет \emph{аналитичность} каждого блока T-матрицы и переносит всю «сложность» геометрии в линейную алгебру \eqref{decomp-eq:gmm-system}, тогда как отдельный численный решатель для коробки разрушил бы единый аналитический каркас и контролируемость погрешности.

\subsection{Протокол верификации}
\label{decomp-ssec:verify}

Верификация многоуровневая: (1) сопоставление с \emph{независимым аналитическим пределом} (редукция к более простой геометрии), (2) сопоставление с \emph{опубликованным численным эталоном}, (3) проверка \emph{солвер-агностических} физических инвариантов. Принятая временна\'я конвенция $e^{-i\omega t}$, исходящая волна $\sim h_n^{(1)}$ (Риккати $\xi_n=\psi_n-i\chi_n$); знак $\operatorname{Im}\varepsilon>0$ отвечает поглощению.

\paragraph{(a) Слоистая сфера: Аден--Керкер, Борен--Хаффман и собственный независимый эталон.}
Импедансно-адмитансная рекурсия по концентрическим оболочкам тождественно совпадает с рекурсией Адена--Керкера для двух сфер и её многослойным обобщением \cite{aden1951}, а в однослойном пределе --- с замкнутыми коэффициентами Ми Борена--Хаффмана \cite{bohren1983}
\begin{align}
a_n &= \frac{m\,\psi_n(mx)\,\psi_n'(x) - \psi_n(x)\,\psi_n'(mx)}{m\,\psi_n(mx)\,\xi_n'(x) - \xi_n(x)\,\psi_n'(mx)},
\label{decomp-eq:mie-a}\\
b_n &= \frac{\psi_n(mx)\,\psi_n'(x) - m\,\psi_n(x)\,\psi_n'(mx)}{\psi_n(mx)\,\xi_n'(x) - m\,\xi_n(x)\,\psi_n'(mx)},
\label{decomp-eq:mie-b}
\end{align}
где $x=k_0 R$ --- размерный параметр, $m=\sqrt{\varepsilon\mu}/\sqrt{\varepsilon_m\mu_m}$ --- относительный показатель преломления, $\psi_n(z)=z\,j_n(z)$ и $\xi_n(z)=z\,h_n^{(1)}(z)$ --- функции Риккати--Бесселя. Решатель проекта дополнительно сверен с \emph{собственным} независимым рядом Ми (модуль \texttt{tests/analytic\_mie.py}), коэффициенты \eqref{decomp-eq:mie-a}--\eqref{decomp-eq:mie-b} которого выведены из замкнутых формул через цилиндрические функции и \emph{не используют} код библиотеки. Достигнута \textbf{машинная точность} ($\sim 10^{-15}$) по $Q_{\mathrm{sca}}$ и $Q_b$ при $x=0{,}5\dots5$, инвариантность к разбиению на $N$ идентичных слоёв, а также корректное поглощение ($Q_{\mathrm{abs}}$ при $\operatorname{Im}\varepsilon>0$) и устойчивый PEC-предел ($|\varepsilon|\to\infty$) --- $10/10$ независимых проверок. Чувствительное к знаку обратное рассеяние $Q_b$ подтвердило корректность конвенции $h_n^{(1)}$.

\paragraph{(b) Сфероид: редукция к сфере при $d\to0$ и эталоны Асано--Ямамото, Фарафонова--Вощинникова, SMARTIES.}
При стягивании фокусного расстояния $d\to0$ вытянутый/сплюснутый сфероид вырождается в сферу, сфероидальные VSWF переходят в сферические, а блочная по $m$ T-матрица --- в диагональную Ми. Это служит первым (аналитическим) контролем. Поля раскладываются по сфероидальным волновым функциям, а граничные условия (непрерывность тангенциальных компонент $\bm{E}$ и $\bm{H}$ на поверхности) дают систему линейных уравнений для коэффициентов разложения \cite{asano1975}; азимутальное число $m$ при этом сохраняется (моды по $m$ расщепляются), тогда как мультипольные порядки внутри фиксированного $m$ связаны, а раздельное TE/TM-описание возможно лишь при осевом ($m=0$) падении вдоль оси симметрии \cite{asano1975,flammer1957}. Сечения и индикатриса сверяются с классическими расчётами рассеяния на сфероиде Асано--Ямамото \cite{asano1975} и с многослойным (конфокальным) обобщением Фарафонова--Вощинникова \cite{farafonov2012,farafonov1996}, а также с независимым кодом сфероидального T-матричного метода (SMARTIES) \cite{mishchenko2002}.

\paragraph{(c) Квазистатический эллипсоид: предел деполяризации и Сихвола--Линделл.}
Для однородного диэлектрического эллипсоида поляризуемость вдоль главной оси $j$
\begin{equation}
\alpha_j = V\,\frac{\varepsilon-\varepsilon_m}{\varepsilon_m + L_j\,(\varepsilon-\varepsilon_m)},
\qquad \sum_{j=x,y,z} L_j = 1,
\label{decomp-eq:depol}
\end{equation}
где $V=\tfrac{4}{3}\pi abc$ --- объём, $L_j$ --- факторы деполяризации \cite{bohren1983}. Контроль: в пределе сферы $L_x=L_y=L_z=\tfrac13$ выражение \eqref{decomp-eq:depol} даёт классическую клаузиус-моссоттиеву поляризуемость $\alpha=3V\,(\varepsilon-\varepsilon_m)/(\varepsilon+2\varepsilon_m)=4\pi a^3\,(\varepsilon-\varepsilon_m)/(\varepsilon+2\varepsilon_m)$. Многослойный (конфокальный) эллипсоид сверяется с $2\times2$-передаточной матрицей факторов деполяризации Сихволы--Линделла \cite{sihvola1990}.

\paragraph{(d) GMM: две сферы против эталона Сюя (Xu).}
Кластерный решатель \eqref{decomp-eq:gmm-system} в конфигурации двух сфер сверяется с опубликованными эталонными сечениями и индикатрисами Сюя \cite{xu1995}; контролируется сходимость по $n_{\max}$ и по порядку многократного рассеяния, а также корректность трансляционных коэффициентов \cite{stein1961,cruzan1962}.

\paragraph{(e) Солвер-агностические инварианты: оптическая теорема и взаимность.}
Независимо от примитива выполняются два физических критерия. \emph{Закон сохранения энергии / оптическая теорема}: сечение экстинкции связано с амплитудой рассеяния вперёд,
\begin{equation}
C_{\mathrm{ext}} = \frac{4\pi}{k}\,\operatorname{Im}\bigl[\hat{\bm{e}}^{*}\!\cdot\bm{S}(\hat{\bm{n}}_{\mathrm{inc}})\bigr],
\qquad C_{\mathrm{ext}} = C_{\mathrm{sca}} + C_{\mathrm{abs}} \ge C_{\mathrm{sca}},
\label{decomp-eq:optical-theorem}
\end{equation}
причём для непоглощающего тела $C_{\mathrm{abs}}=0$ \cite{bohren1983,mishchenko2002}. \emph{Взаимность} налагает на матрицу рассеяния (и на T-матрицу) симметрию относительно перестановки направлений падения и наблюдения \cite{mishchenko2002,waterman1971}. Оба критерия не зависят от конкретного решателя и потому служат сквозной проверкой связки «примитив $\to$ T-матрица $\to$ GMM»; нарушение \eqref{decomp-eq:optical-theorem} немедленно сигнализирует об ошибке в знаках/конвенции или о незаконном сопряжении амплитуд.

\subsection{Сводная таблица: решатель $\leftrightarrow$ эталон}
\label{decomp-ssec:verif-table}

\begin{table}[htbp]
\centering
\caption{Соответствие решателей примитивов их верификационным эталонам и достигнутой точности.}
\label{decomp-tab:verif-map}
\small
\begin{tabular}{@{}p{3.0cm}p{3.0cm}p{4.2cm}p{2.4cm}@{}}
\toprule
\textbf{Решатель} & \textbf{Аналитический предел} & \textbf{Опубликованный/независимый эталон} & \textbf{Точность} \\
\midrule
Слоистая сфера &
Однослойный Ми; $N$ идентичных слоёв $\equiv$ 1 слой &
Аден--Керкер \cite{aden1951}; Борен--Хаффман \cite{bohren1983}; собственный ряд Ми \texttt{analytic\_mie.py} &
$\sim10^{-15}$ ($10/10$) \\[2pt]
Конфокальный сфероид &
$d\to0\Rightarrow$ сфера (Ми) &
Асано--Ямамото \cite{asano1975}; Фарафонов--Вощинников \cite{farafonov2012}; SMARTIES \cite{mishchenko2002} &
эталонная \\[2pt]
Квазистат. эллипсоид &
$L_j=\tfrac13\Rightarrow$ сфера (Клаузиус--Моссотти) &
Сихвола--Линделл \cite{sihvola1990}; факторы $L_j$ \cite{bohren1983} &
аналитическая \\[2pt]
GMM (кластер) &
$P=1\Rightarrow$ одиночный примитив &
Две сферы, Сюй \cite{xu1995}; трансл. теоремы \cite{stein1961,cruzan1962} &
эталонная \\[2pt]
\textbf{Любой} (сквозной) &
$C_{\mathrm{abs}}\!=\!0$ для непоглощающего тела &
Оптическая теорема \& взаимность \cite{mishchenko2002,bohren1983} &
инвариант \\
\bottomrule
\end{tabular}
\end{table}

Совокупность проверок \emph{(a)--(e)} образует иерархию доверия: машинно-точная верификация ядра-сферы \cite{aden1951,bohren1983}, эталонная сверка сфероида \cite{asano1975,farafonov2012} и эллипсоида \cite{sihvola1990}, кросс-валидация кластера \cite{xu1995} и, поверх всего, не зависящие от реализации законы сохранения \cite{mishchenko2002}. Именно эта многоуровневость, а не отдельный «золотой» тест, обеспечивает корректность всего аналитического каркаса при переходе от единичного примитива к сложной составной геометрии.

\begin{thebibliography}{99}
\bibitem{eisenhart1934} L.~P.~Eisenhart, \emph{Separable systems of St\"ackel}, Ann. Math. \textbf{35}, 284--305 (1934).
\bibitem{stratton1941} J.~A.~Stratton, \emph{Electromagnetic Theory} (McGraw-Hill, New York, 1941).
\bibitem{aden1951} A.~L.~Aden, M.~Kerker, \emph{Scattering of electromagnetic waves from two concentric spheres}, J. Appl. Phys. \textbf{22}, 1242--1246 (1951).
\bibitem{flammer1957} C.~Flammer, \emph{Spheroidal Wave Functions} (Stanford Univ. Press, 1957).
\bibitem{stein1961} S.~Stein, \emph{Addition theorems for spherical wave functions}, Q. Appl. Math. \textbf{19}, 15--24 (1961).
\bibitem{cruzan1962} O.~R.~Cruzan, \emph{Translational addition theorems for spherical vector wave functions}, Q. Appl. Math. \textbf{20}, 33--40 (1962).
\bibitem{wittmann1988} R.~C.~Wittmann, \emph{Spherical wave operators and the translation formulas}, IEEE Trans. Antennas Propag. \textbf{36}, 1078--1087 (1988).
\bibitem{waterman1971} P.~C.~Waterman, \emph{Symmetry, unitarity, and geometry in electromagnetic scattering}, Phys. Rev. D \textbf{3}, 825--839 (1971).
\bibitem{asano1975} S.~Asano, G.~Yamamoto, \emph{Light scattering by a spheroidal particle}, Appl. Opt. \textbf{14}, 29--49 (1975).
\bibitem{bohren1983} C.~F.~Bohren, D.~R.~Huffman, \emph{Absorption and Scattering of Light by Small Particles} (Wiley, New York, 1983).
\bibitem{sihvola1990} A.~H.~Sihvola, I.~V.~Lindell, \emph{Polarizability and effective permittivity of layered and continuously inhomogeneous dielectric ellipsoids}, J. Electromagn. Waves Appl. \textbf{4}, 1--26 (1990).
\bibitem{xu1995} Y.-L.~Xu, \emph{Electromagnetic scattering by an aggregate of spheres}, Appl. Opt. \textbf{34}, 4573--4588 (1995).
\bibitem{farafonov1996} V.~G.~Farafonov, N.~V.~Voshchinnikov, V.~V.~Somsikov, \emph{Light scattering by a core-mantle spheroidal particle}, Appl. Opt. \textbf{35}, 5412--5426 (1996).
\bibitem{mishchenko2002} M.~I.~Mishchenko, L.~D.~Travis, A.~A.~Lacis, \emph{Scattering, Absorption, and Emission of Light by Small Particles} (Cambridge Univ. Press, 2002).
\bibitem{farafonov2012} V.~G.~Farafonov, N.~V.~Voshchinnikov, \emph{Light scattering by a multilayered spheroidal particle}, Appl. Opt. \textbf{51}, 1586--1597 (2012); arXiv:1110.6604.
\bibitem{vanburen2020} A.~L.~Van Buren, J.~E.~Boisvert, \emph{Accurate calculation of oblate spheroidal wave functions with complex argument}, arXiv:2009.01618 (2020).
\bibitem{dlmf} \emph{NIST Digital Library of Mathematical Functions}, \url{https://dlmf.nist.gov/}.
\bibitem{bowman1987} J.~J.~Bowman, T.~B.~A.~Senior, P.~L.~E.~Uslenghi (eds.), \emph{Electromagnetic and Acoustic Scattering by Simple Shapes} (Hemisphere, New York, 1987).
\bibitem{klinkenbusch2007} L.~Klinkenbusch, \emph{Electromagnetic scattering by semi-infinite circular and elliptic cones}, Radio Sci. \textbf{42}, RS6S10 (2007).
\bibitem{daylis2017} С.~Дайлис, С.~Н.~Шабунин, \emph{Функции Грина многослойных цилиндрических структур в задачах излучения, распространения и дифракции электромагнитных волн}, Ural Radio Engineering Journal \textbf{1}(1), 91--117 (2017); DOI 10.15826/urej.2017.1.1.005.
\bibitem{mitelman2013} Ю.~Е.~Мительман, С.~Н.~Шабунин, \emph{Электродинамика многослойных цилиндрических направляющих систем} (LAP Lambert Academic Publishing, 2013).
\bibitem{wait1955} J.~R.~Wait, \emph{Scattering of a plane wave from a circular dielectric cylinder at oblique incidence}, Can. J. Phys. \textbf{33}, 189--195 (1955).
\bibitem{kavaklioglu2008} Ö.~Kavaklıoğlu, B.~Schneider, \emph{On the Fourier--Bessel multiple scattering coefficients of an infinite grating of dielectric circular cylinders at oblique incidence} (Wait's isolated-cylinder coefficients), arXiv:0711.0362 (2008).
\end{thebibliography}
\end{document}
